% !TEX program = pdflatex
% Standalone categorical development of open harmonic chains and divisible
% strings, companion to ../CDLM.tex.  CDLM results are cited by the numbering
% of the 2026-07 draft; each such cite carries a comment recording the CDLM
% label, so numbers can be refreshed from CDLM.aux if the draft renumbers.
\documentclass[11pt,oneside,article]{memoir}

\linespread{1.12}\selectfont
\settrims{0pt}{0pt}
\settypeblocksize{*}{35.4pc}{*}
\setlrmargins{*}{*}{1}
\setulmarginsandblock{0.85in}{0.85in}{*}
\setheadfoot{\onelineskip}{2.66\onelineskip}
\setheaderspaces{*}{1.5\onelineskip}{*}
\checkandfixthelayout[lines]

\usepackage[T1]{fontenc}
\usepackage[utf8]{inputenc}
\usepackage{amsmath,amssymb,mathtools,amsthm}
\usepackage{lmodern}
\usepackage{amscd}
\usepackage{tikz-cd}
\usepackage{xcolor}
\usepackage[inline]{enumitem}
\usepackage[bookmarks=true,colorlinks=true,linkcolor=blue!50!black,
  citecolor=orange!50!black,urlcolor=orange!50!black]{hyperref}
\usepackage[capitalize]{cleveref}

\setlist{nosep}
\setlength{\emergencystretch}{3em}
\setsecnumdepth{subsection}
\settocdepth{subsection}
\renewcommand{\thesection}{\arabic{section}}
\renewcommand{\thesubsection}{\thesection.\arabic{subsection}}

\theoremstyle{definition}
\newtheorem{definitionx}{Definition}[section]
\newtheorem{examplex}[definitionx]{Example}
\theoremstyle{plain}
\newtheorem{theorem}[definitionx]{Theorem}
\newtheorem{proposition}[definitionx]{Proposition}
\newtheorem{lemma}[definitionx]{Lemma}
\newtheorem{corollary}[definitionx]{Corollary}
\theoremstyle{remark}
\newtheorem{remarkx}[definitionx]{Remark}

% End-of-definition lozenge, as in CDLM.
\newenvironment{definition}
  {\pushQED{\qed}\renewcommand{\qedsymbol}{$\lozenge$}\definitionx}
  {\popQED\enddefinitionx}
\newenvironment{example}
  {\pushQED{\qed}\renewcommand{\qedsymbol}{$\lozenge$}\examplex}
  {\popQED\endexamplex}
\newenvironment{remark}
  {\pushQED{\qed}\renewcommand{\qedsymbol}{$\lozenge$}\remarkx}
  {\popQED\endremarkx}
\crefname{definitionx}{Definition}{Definitions}
\crefname{examplex}{Example}{Examples}
\crefname{remarkx}{Remark}{Remarks}

% ---------- basic macros ----------
\newcommand{\R}{\mathbb R}
\newcommand{\N}{\mathbb N}
\newcommand{\Z}{\mathbb Z}
\newcommand{\id}{\mathrm{id}}
\newcommand{\blank}{\mathord{-}}
\newcommand{\To}[1]{\xrightarrow{#1}}
\newcommand{\inpt}[1]{#1^{-}}% input part (as in CDLM)
\newcommand{\outp}[1]{#1^{+}}% output part (as in CDLM)
\newcommand{\Cat}[1]{\mathbf{#1}}
\newcommand{\cat}[1]{\mathcal{#1}}
\newcommand{\Fn}[1]{\textnormal{\textsf{#1}}}% a named functor
\newcommand{\fun}[1]{\mathrm{#1}}% a named function or operation
\newcommand{\inc}{\fun{inc}}
\newcommand{\rv}[1]{\mathbf{#1}}% a reactive Hilbert space (V,\sharpR_V)
\newcommand{\ltens}{\boxtimes}% lens/interface tensor
\newcommand{\phase}{\mathrm{phase}}
\newcommand{\ParaCat}{\mathbb{P}\Cat{ara}}
\newcommand{\para}[2]{\ParaCat_{#1}(#2)}
\newcommand{\cotof}[1]{\cot(#1)}

% ---------- term tracking (from CDLM) ----------
% \trackTerm{\xyz}{symbol} defines \xyz to typeset symbol; every use of \xyz
% hyperlinks (in black) back to the definition site \defineTerm{xyz}.  The
% tracking key is derived from the macro name itself, so key = macro holds by
% construction; \defineTerm with any other key is a compile error.  The
% starred form \trackTerm* overwrites an existing command (e.g. \cot).
\makeatletter
\NewDocumentCommand{\termraw}{m}{%
  \@ifundefined{termraw@#1}%
    {\PackageError{termraw}{Term '#1' not registered}{Register it with \string\trackTerm first.}}%
    {\@nameuse{termraw@#1}}%
}
\DeclareRobustCommand{\termref}[2]{{\hypersetup{linkcolor=black}\hyperlink{term.#1}{#2}}}
\NewDocumentCommand{\trackTerm}{s m m}{%
  \edef\trackTerm@key{\expandafter\@gobble\string#2}%
  \expandafter\newcommand\csname termraw@\trackTerm@key\endcsname{{#3}}%
  \IfBooleanTF{#1}{\let#2\@undefined}{}%
  \expandafter\trackTerm@def\expandafter{\trackTerm@key}{#2}%
}
\NewDocumentCommand{\trackTerm@def}{m m}{%
  \newcommand{#2}{\termref{#1}{\termraw{#1}}}%
}
\NewDocumentCommand{\defineTerm}{m}{\hypertarget{term.#1}{\termraw{#1}}}
\makeatother

\trackTerm{\PreHilb}{\Cat{PreHilb}}
\trackTerm{\Hilb}{\Cat{Hilb}}
\trackTerm{\hsm}{\Cat{Hilb}_{\mathrm{sm}}}
\trackTerm{\rhilb}{\Cat{RHilb}}
\trackTerm{\smsetiso}{\Cat{Set}_{\cong}}
\trackTerm{\poly}{\Cat{Poly}}
\trackTerm{\yon}{\mathsf{y}}
\trackTerm{\pc}{\mathbb{P}\Cat{C}}
\trackTerm{\harr}{\mathbb{A}\Cat{rr}_{\mathrm{Hb}}}
\trackTerm{\Dyn}{\Cat{Dyn}}
\trackTerm{\Int}{\Cat{Int}}
\trackTerm{\DelayCat}{\Cat{Delay}}
\trackTerm{\Store}{\Fn{Store}}
\trackTerm{\Hb}{\mathrm{Hb}}% the Hilbert adaptive-arrangement datum
\trackTerm{\interph}{\mathrm{hb}}% the Hilbert polynomial interpretation datum
\trackTerm{\thetacomp}{\theta}% effect comparison of the interpretation datum
\trackTerm{\potd}{\fun{d}}% potential algebra map cot(R) -> y
\trackTerm{\chiE}{\chi}% Euler step
\trackTerm{\nuro}{\nu}% phase readout
\trackTerm{\dynphase}{\fun{dyn}_{\phase}}% phase dynamics
\trackTerm{\step}{\operatorname{step}}% queue / delay-line step
\trackTerm{\EZ}{E_Z}% completed characteristic state
\trackTerm{\CZ}{C_Z}% cell packet
\trackTerm{\riesz}{r}% Riesz map
\trackTerm{\sharpR}{\sharp}% reaction sharp
\trackTerm{\sharpS}{\sharp^{\mathrm{s}}}% canonical symplectic sharp
\trackTerm{\Open}{\Cat{Open}}% semicategory of open harmonic systems (1-categorical: \Cat, not \mathbb)
\trackTerm{\Lin}{\mathbb{L}\Cat{in}}% semicategory of ported machines
\trackTerm{\Ufun}{\mathsf{U}}% embedding of ported machines into PC
\trackTerm{\Phiopen}{\Phi_{\mathrm{open}}}% open linear semantics
\trackTerm{\seam}{\fun{seam}}% seam wiring
\trackTerm{\Qu}{Q}% queue state space Q_N(B)
\trackTerm{\Qterm}{\mathbb{Q}}% terminal delay realization
\trackTerm{\Fl}{\Fn{Fl}}% covector fields on a Hilbert space
\trackTerm{\blk}{\fun{blk}}% block steps as a natural transformation
\trackTerm{\headf}{\fun{head}}% completed-step head
\trackTerm{\nextf}{\fun{next}}% completed-step next
\trackTerm{\Zport}{P_Z}% the Z-port
\trackTerm{\cellm}{\fun{cell}}% the cell
\trackTerm{\Chain}{\fun{chain}}% chains = powers of the cell
\trackTerm{\Tm}{\mathsf{T}}% the chain machine
\trackTerm{\Sm}{\mathsf{S}}% the string
\trackTerm{\Phip}{\Phi'}% polynomial interpretation functor
\trackTerm*{\cot}{\Fn{cot}}
\trackTerm{\Phiphase}{\Phi_{\phase}}

\title{Open Harmonic Chains and Divisible Strings}
\author{}
\date{}

\begin{document}
\maketitle

\begin{abstract}
The CDLM framework of adaptive arrangements is parametric in an
adaptive-arrangement datum.  We assemble a \emph{Hilbert arrangement
datum} $\Hb$---spaces are real Hilbert spaces with smooth maps,
parameters are reactive Hilbert spaces, potentials accumulate in
$\R$---and obtain the operad $\harr$ of Hilbert arrangements and its
phase semantics $\Phiphase\colon\harr\to\pc$.  On open harmonic-linear
systems the semantics refines: an interface becomes a \emph{port}, a
Hilbert space equipped with an impedance form; a system becomes a
bounded linear \emph{ported machine}; serial gluing of systems is
substitution into a seam wiring; and the refinement
$\Phiopen\colon\Open\to\Lin$ is a semifunctor carrying gluing to
composition of machines.

One open system generates everything: the \emph{cell}, a unit mass
bound to its input port at impedance $Z$.  Chains are its powers,
$\Chain_{K+K'}=\Chain_{K'}\circ\Chain_K$, and the semantics turns this
equation into serial composition of chain machines.  Each chain
machine carries an exact characteristic flow: in wave variables its
driven trajectories transport packets, which enter at one port and
exit at the other exactly $K$ ticks later, and from the first transit
onward the wave chart coincides with the state of the standard
two-queue string.  Strings subdivide, by packet doubling and time
change along multiplication by two on discrete time, and the Hilbert
completion of the dyadic tower is the weighted space $L^2([0,d),Z\,dx)$
in each travel direction, on which every step is cut-and-translate:
the \emph{divisible string}, a strong semigroupal representation of
durations in which concatenation of intervals is serial composition.
The infinite-resolution chain is thus the completion of its finite
open chains, and translation of the two completed fields is equivalent
to the weak wave equation.
\end{abstract}

\tableofcontents*

\section{Introduction}
\label{sec.introduction}

This paper constructs the infinite-resolution harmonic chain as the
completion of its finite open chains, with a categorical statement at
every step.  This introduction lists the statements.

\paragraph{Syntax.}
The framework of \cite{cdlm} is parametric in an adaptive-arrangement
datum; \cref{sec.hilbert_datum} assembles the Hilbert arrangement
datum $\Hb=(\hsm,\rhilb,\inc,\R)$---spaces are real Hilbert spaces
with smooth maps, parameters are reactive Hilbert spaces, potentials
accumulate in $\R$---and obtains the operad $\harr$ of Hilbert
arrangements.  The boundary datum of an open harmonic system is a
\emph{port}, a Hilbert space with an impedance form
(\cref{def.port}).  Open systems form a semicategory $\Open$
(\cref{def.open}): composition is substitution into a seam wiring, so
associativity is the substitution law of the operad applied to a
single equality of wirings (\cref{lem.seams_associate}), and direct
sum makes $\Open$ symmetric semigroupal.  There are no identities,
because output maps of arrangements read only the parameter, and the
absence descends through the semantics to the nonexistence of strings
of duration zero.

\paragraph{Semantics.}
The interpretation datum and the phase integrator of
\cref{sec.hilbert_semantics} are assembled from monoidal functors and
monoidal natural transformations, and produce the normal operad
functor $\Phiphase\colon\harr\to\pc$ into the $2$-category of
polynomial coalgebras (\cref{thm.phase_functor}).  On open systems
the semantics refines: ported machines form a semicategory $\Lin$
(\cref{def.ported_machine}), embedded faithfully in $\pc$ by affine
covector fields (\cref{def.ufun}), and the open linear semantics
$\Phiopen\colon\Open\to\Lin$ is the unique semifunctor compatible
with $\Phiphase$ (\cref{thm.open_semantics}).

\paragraph{Machines across ports.}
Machines at different ports compare along port embeddings, the
isometries intertwining the impedance operators
(\cref{def.port_embedding}).  Morphisms over port embeddings carry
values along the embeddings and covectors along their Riesz
conjugates (\cref{def.machine_over}); they compose horizontally
(\cref{lem.horizontal_composition}); and towers of them have colimits
computed by completed unions (\cref{prop.machine_colimits})---the
machine-level image of the completion adjunction between pre-Hilbert
and Hilbert spaces (\cref{sec.background_hilbert}).

\paragraph{The generator.}
One morphism generates the mechanics: the cell, a unit mass bound to
its input port at impedance $Z$ (\cref{def.cell}).  Chains are its
powers, and $K\mapsto\Chain_K$ is a homomorphism of semigroups
$(\Z_{\geq1},+)\to\Open(\Zport,\Zport)$
(\cref{cor.chain_semigroup}), carried by $\Phiopen$ to the strictly
composing chain machines (\cref{cor.chain_machines_compose}).

\paragraph{Trajectories.}
A driven trajectory of a ported machine is a $2$-cell from the
successor system $(\N,\mathrm{succ})$ into the closure of the machine
by a stream environment
(\cref{def.stream_environment,prop.trajectories_are_elements}): the
trajectories are the $(\N,\mathrm{succ})$-shaped elements of a closed
system in $\pc$.  The characteristic flow theorem
(\cref{thm.characteristic_flow}) is proved by direct computation with
the recurrences of the chain machine: driven chain trajectories
transport wave packets, which enter at one port and exit at the other
exactly $K$ ticks later.  Its output is recaptured categorically at
once: from the first transit onward, the wave chart of the chain
trajectory is the state of the standard two-queue string
(\cref{cor.chain_string}).

\paragraph{Queues, strings, and resolution.}
The queue is terminal among bounded linear realizations of exact
delay (\cref{thm.queue_representation}).  Strings compose:
$K\mapsto\Sm_B(K)$ is a strong semigroupal representation of
$(\Z_{\geq1},+)$ in the endomorphism category of a zero-impedance
port (\cref{prop.string_composition}), and the string is a functor of
its packet alphabet that preserves tower colimits
(\cref{prop.string_functorial}).  Subdivision is time change along
multiplication by two on discrete time
(\cref{thm.queue_subdivision,cor.dyadic_subdivision}).  The block
steps of consecutive resolutions form a natural transformation
between tower functors, and the completed step is its colimit
(\cref{lem.block_naturality,thm.step_colimit}); the completed state
is the colimit in Hilbert spaces and contractions, realized as
$L^2\bigl([0,d),Z\,dx\bigr)$ (\cref{thm.dyadic_completion}).

\paragraph{Divisibility.}
Weighted interval spaces are a strong algebra over the ordered
interval operad (\cref{thm.characteristic_interval_algebra}).  The
divisible string $d\mapsto\Sm_Z(d;a)$ is a strong semigroupal
representation of $(\R_{\geq a},+)$ in which concatenation of
intervals is serial composition (\cref{thm.divisibility}); it is the
colimit of the finite strings over the block alphabets, and the
identifications carry the finite concatenation $2$-cells to the
commensurable concatenations (\cref{thm.string_colimit}).  The main
theorem (\cref{thm.main}) assembles these: gluing of cells, serial
composition of machines, concatenation of queues, and concatenation
of intervals are one composition, and the infinite-resolution chain
is the completion of its finite open chains.

\paragraph{Where analysis enters.}
Analysis enters in three places.  The Riesz dagger structure of
$\Hilb$ (\cref{sec.background_hilbert}) feeds the covector legs of
every backward pass.  The completion adjunction between pre-Hilbert
and Hilbert spaces, packaged once as \cref{lem.isometric_colimits},
is consumed by every colimit in the paper.  And the local $L^2$
calculus of \cref{sec.waves} converts transport of the two completed
fields into weak solutions of the wave equation, by the
distributional Poincar\'e lemma (\cref{thm.characteristic_wave}), and
the energy space of a solution into the state space of the divisible
string (\cref{cor.energy_coordinates}).

\section{Assumed background}
\label{sec.background}

This paper is a companion to \cite{cdlm}, whose framework it instantiates;
we cite its results by number.  This section records everything else we
assume.  \Cref{sec.background_categories} fixes the categorical background
and recalls polynomial functors and their coalgebras;
\cref{sec.background_hilbert} fixes the Hilbert-space background.  Terms
introduced in the body of the paper are marked at their definition sites,
and every later use of the corresponding symbol hyperlinks back to that
site.

\subsection{Categorical background}
\label{sec.background_categories}

We assume familiarity with symmetric monoidal categories; with monoids and
comonoids in them; with lax, strong, and normal (unit-preserving) monoidal
functors and monoidal natural transformations; with monads, monoidal monads,
and adjunctions; and with the writer monad $R\otimes\blank$ of a commutative
monoid object $(R,0,+)$ \cite[Lem.~3.2.1]{cdlm}.%lem.monoid_to_monad
We write $\defineTerm{smsetiso}$ for the category of sets and bijections.

\paragraph{Operads.}
By an \emph{operad} we mean a symmetric colored operad, i.e.\ a symmetric
multicategory \cite{leinster2004higher}: a class of \emph{colors}; for every
finite list of colors $c_1,\ldots,c_K$ and every color $c$ a set of
\emph{operations} $\cat O(c_1,\ldots,c_K;c)$; identities; substitution; and
symmetric-group actions, subject to associativity, unit, and equivariance
laws.  Every symmetric monoidal category $(\cat C,I,\otimes)$ has an
underlying operad with
$\cat O(c_1,\ldots,c_K;c)\coloneqq\cat C(c_1\otimes\cdots\otimes c_K,c)$,
and every lax symmetric monoidal functor induces an \emph{operad functor}
between the underlying operads.

For a nonsymmetric colored operad $\cat O$ and a symmetric monoidal category
$(\cat C,I,\otimes)$, an \emph{$\cat O$-algebra in $\cat C$} assigns an
object $A(c):\cat C$ to each color and a map
$A(\varphi)\colon A(c_1)\otimes\cdots\otimes A(c_K)\to A(c)$ to each
operation $\varphi\in\cat O(c_1,\ldots,c_K;c)$, respecting identities and
substitution.  We denote such an algebra as a map $A\colon\cat O\to\cat C$.
We call the maps $A(\varphi)$ the \emph{structure maps} of $A$, and we call
$A$ \emph{strong} if every structure map is an isomorphism.

\paragraph{Semicategories and semigroupal categories.}
A \emph{semicategory} is a category without the assumption of identity
morphisms: objects, hom-sets, and an associative composition.  A
\emph{semifunctor} between semicategories assigns objects to objects and
morphisms to morphisms, preserving composition.  A \emph{semigroupal
category} is a category with a tensor product and an associator satisfying
the pentagon, with no unit assumed---the nonunital version of a monoidal
category.  A \emph{strong semigroupal functor} is a functor equipped with
natural isomorphisms $F(c)\otimes F(c')\cong F(c\otimes c')$ satisfying the
associativity coherence.  A commutative semigroup, such as
$(\R_{\geq a},+)$, is a discrete semigroupal category; a \emph{strong
semigroupal representation} of it in a semigroupal category $\cat D$ is a
strong semigroupal functor into $\cat D$.

\paragraph{Polynomial functors.}
We use the category $\defineTerm{poly}$ of polynomial functors; see
\cite{niu2025polynomial} for a full treatment and \cite[\S2]{cdlm} for the
fragment used here.  A polynomial is a formal coproduct of representables
\[
p=\sum_{i\in p(1)}\yon^{p[i]},
\]
where the elements of the set $p(1)$ are the \emph{positions} of $p$ and the
elements of each set $p[i]$ are the \emph{directions} at $i$.  A morphism
$p\to p'$ is a function forward on positions together with, at each
$i\in p(1)$, a function backward on directions
$p'[i']\to p[i]$.  We write $\defineTerm{yon}$ for the identity polynomial
(one position, one direction), so $\yon$ is the monoidal unit for the
Dirichlet product $\otimes$, which multiplies position sets and direction
sets.  The Dirichlet product is closed; we write $[p,p']$ for the internal
hom.  For a set $S$, we write $\defineTerm{Store}(S)\coloneqq S\yon^S$.  A
\emph{$p$-coalgebra} is a set $S$ with a function $f\colon S\to p(S)$;
equivalently \cite[Prop.~2.2.3]{cdlm},%prop.coalg_as_poly_map
a polynomial map $\Store(S)\to p$.

\paragraph{Polynomial coalgebras.}
We use the $2$-category $\defineTerm{pc}$ of \cite[\S2.2.4]{cdlm}:%sec.org
its objects are polynomials; its hom-category $\pc(p,p')$ is the category of
$[p,p']$-coalgebras and coalgebra morphisms; composition and the monoidal
structure (with unit $\yon$) are by Dirichlet product of coalgebras.  We
refer to objects of $\pc(p,p')$ as \emph{$1$-cells} and to their morphisms
as \emph{$2$-cells}.  Since $[\yon,\yon]\cong\yon$, a $1$-cell
$\yon\to\yon$ is a $\yon$-coalgebra, i.e.\ a set $S$ with a function
$S\to S$---a discrete dynamical system---and a $2$-cell between two such is
a function commuting with the two endofunctions.

\begin{lemma}[Stateless 1-cells]\label{lem.stateless}
For polynomials $p,p'$, evaluation gives
$[p,p'](1)\cong\poly(p,p')$, so every polynomial map
$f\colon p\to p'$ determines a $1$-cell in $\pc(p,p')$ with a single
state, its \emph{stateless} $1$-cell.  The assignment preserves
identities, hom-composition, and $\otimes$; we elide it, writing $f$
for the stateless $1$-cell of $f$.
\end{lemma}

\begin{proof}
A coalgebra with carrier $1$ is a point of $[p,p'](1)$, and
$[p,p'](1)\cong\poly(p,p')$ \cite{niu2025polynomial}.  On single-state
carriers, hom-composition and $\otimes$ of coalgebras restrict to
composition and $\otimes$ of the selected polynomial maps, by direct
unwinding of the Dirichlet product of coalgebras.
\end{proof}

\paragraph{The CDLM framework.}
The framework of \cite[Ch.~4]{cdlm} is parametric in an
\emph{adaptive-arrangement datum} $D=(\cat M,\cat Q,J,\Fn R)$: a cartesian
monoidal category $\cat M$ of \emph{spaces}, a symmetric monoidal category
$\cat Q$ of \emph{parameters}, a strong monoidal functor
$J\colon\cat Q\to\cat M$, and a monoidal monad $\Fn R$ on $\cat M$, the
\emph{potentials monad} \cite[Def.~4.1.1]{cdlm}.%def.rewiring_datum
Out of $D$ the framework builds an operad $\mathbb{A}\Cat{rr}_D$ of
\emph{adaptive arrangements}; a \emph{polynomial interpretation datum} for
$D$ \cite[Def.~4.2.1]{cdlm}%def.polynomial_interpretation
then yields a lax symmetric monoidal functor into parameterized polynomials
\cite[Thm.~4.2.2]{cdlm},%thm.poly_interpretation
and an \emph{integrator} \cite[Def.~4.3.4]{cdlm}%def.integrator
turns the result into an operad functor
$\mathbb{A}\Cat{rr}_D\to\pc$ \cite[Thm.~4.3.6]{cdlm}.%thm.dynamics_functor
We recall the relevant definitions at their points of use in
\cref{sec.hilbert_datum,sec.hilbert_semantics}, where we instantiate each of
them in Hilbert spaces.

\subsection{Hilbert-space background}
\label{sec.background_hilbert}

All vector spaces are real.  Standard references for this material are
\cite{conway1990functional} for Hilbert-space theory,
\cite{cartan1971differential} for differential calculus on Banach spaces,
and \cite{brezis2011functional} for Sobolev spaces and distributions.

\paragraph{Pre-Hilbert and Hilbert spaces.}
A \emph{pre-Hilbert space} is a vector space with an inner product
$\langle\blank,\blank\rangle$; a \emph{Hilbert space} is a pre-Hilbert space
that is complete for the induced norm.  A linear map $f\colon V\to W$ is
\emph{bounded} if $\lVert f\rVert\coloneqq\sup_{\lVert v\rVert\leq1}\lVert
f(v)\rVert_W<\infty$; we call $\lVert f\rVert$ the \emph{operator norm} of
$f$.  We write $\defineTerm{PreHilb}$ for the category of
pre-Hilbert spaces and bounded linear maps, and $\defineTerm{Hilb}$ for its
full subcategory of Hilbert spaces.  A bounded linear map is an
\emph{isometry} if it preserves norms, and \emph{unitary} if it is a
surjective isometry.  In any category, a morphism is a \emph{split
epimorphism} if it admits a section.  For $c\in\R_{>0}$, we write $V[c]$
for the same vector space with the \emph{rescaled} inner product
$c\langle\blank,\blank\rangle$; it is again a pre-Hilbert space, complete
if $V$ is.

\paragraph{Completion.}
The inclusion $\Hilb\hookrightarrow\PreHilb$ has a left adjoint, the
\emph{Hilbert completion} $V\mapsto\widehat V$: the unit $V\to\widehat V$ is
an isometry with dense image, and every bounded linear map $f\colon V\to Y$
with $Y$ complete extends uniquely to $\widehat f\colon\widehat V\to Y$ with
$\lVert\widehat f\rVert=\lVert f\rVert$.

\paragraph{Increasing unions.}
Given pre-Hilbert spaces $H_n$ and isometric embeddings
$\rho_n\colon H_n\hookrightarrow H_{n+1}$ for $n\in\N$, the \emph{increasing
union} $\bigcup_n H_n$ is the union of the images in the colimit of the
underlying vector spaces, a pre-Hilbert space in which each $H_n$ embeds
isometrically and every element lies in some $H_n$.

\begin{lemma}[Isometric colimits]\label{lem.isometric_colimits}
Let $(H_n,\rho_n)_{n\in\N}$ be as above and let $Y\in\Hilb$.  Restriction
along the embeddings $H_n\hookrightarrow\widehat{\bigcup_nH_n}$ is a
bijection
\begin{equation}\label{eq.completion_universal_property}
\Hilb\Bigl(\widehat{\textstyle\bigcup_nH_n},\ Y\Bigr)
\cong
\left\{
(f_n)_{n\in\N}
\ \middle|\
\begin{array}{l}
f_n\colon H_n\to Y\text{ is bounded},\\
f_{n+1}\rho_n=f_n,\
\sup_n\lVert f_n\rVert<\infty
\end{array}
\right\}.
\end{equation}
\end{lemma}

\begin{proof}
Restriction sends a bounded $f$ to a compatible family with
$\lVert f_n\rVert\leq\lVert f\rVert$.  It is injective because
$\bigcup_nH_n$ is dense in the completion.  For surjectivity, let $(f_n)$
be a compatible family with $s\coloneqq\sup_n\lVert f_n\rVert<\infty$.
Every element $x\in\bigcup_nH_n$ lies in some $H_n$; setting
$F(x)\coloneqq f_n(x)$ is well defined by compatibility, linear, and
bounded with $\lVert F\rVert\leq s$.  By the completion adjunction, $F$
extends uniquely to a bounded map on $\widehat{\bigcup_nH_n}$, which
restricts to the given family.
\end{proof}

Restricted to contractions, \cref{eq.completion_universal_property} says
exactly that $\widehat{\bigcup_nH_n}$ is the colimit of $(H_n,\rho_n)$ in
the category of Hilbert spaces and linear contractions.

\paragraph{Duals, Riesz representation, adjoints.}
The \emph{dual} $V^*$ of a pre-Hilbert space $V$ is the space of bounded
linear functionals with the operator norm.  For $V$ a Hilbert space, $V^*$
is a Hilbert space, and the map $V\to V^*$,
$v\mapsto\langle v,\blank\rangle$, is a unitary isomorphism---this is the
Riesz representation theorem \cite{conway1990functional}.  We define the
\emph{Riesz map} $\defineTerm{riesz}_V\colon V^*\To{\cong}V$ to be its inverse.
Consequently every Hilbert space is \emph{reflexive}: the evaluation map
$\iota_V\colon V\to V^{**}$, $\iota_V(v)(\xi)\coloneqq\xi(v)$, is a unitary
isomorphism, and we use it silently to identify $V^{**}$ with $V$.  A
bounded linear map $f\colon V\to W$ has a \emph{transpose}
$f^*\colon W^*\to V^*$, $f^*(\eta)\coloneqq\eta\circ f$, with
$\lVert f^*\rVert=\lVert f\rVert$, $(g\circ f)^*=f^*\circ g^*$, and
$(f^{-1})^*=(f^*)^{-1}$ when $f$ is invertible.  Under reflexivity,
$f^{**}=f$.

For Hilbert spaces, conjugating the transpose by the Riesz maps gives the
\emph{adjoint}
\[
f^\dagger\coloneqq \riesz_V\circ f^*\circ\riesz_W^{-1}\colon W\to V,
\qquad
\langle f(v),w\rangle_W=\langle v,f^\dagger(w)\rangle_V,
\]
and $(\blank)^\dagger$ is a \emph{dagger structure} on $\Hilb$
\cite{heunen2019categories}: an identity-on-objects contravariant
involution, i.e.
\[
\id_V^\dagger=\id_V,
\qquad
(g\circ f)^\dagger=f^\dagger\circ g^\dagger,
\qquad
f^{\dagger\dagger}=f.
\]
Rearranging the definition of the adjoint gives the \emph{Riesz
identities}: for bounded $f\colon V\to W$,
\begin{equation}\label{eq.riesz_identities}
\riesz_V\circ f^*=f^\dagger\circ\riesz_W
\qquad\text{and}\qquad
\riesz_W\circ(f^\dagger)^*=f\circ\riesz_V,
\end{equation}
the second by applying the first to $f^\dagger$ in place of $f$.
In dagger terms the isometries and unitaries of the previous paragraph are
the standard notions: $f$ is an isometry if and only if
$f^\dagger\circ f=\id$ (by polarization), and unitary if and only if it is
a \emph{dagger isomorphism}, $f^\dagger\circ f=\id$ and
$f\circ f^\dagger=\id$.

\paragraph{Direct sums.}
The direct sum $V\oplus W$ of pre-Hilbert spaces carries the inner product
$\langle(v,w),(v',w')\rangle\coloneqq\langle v,v'\rangle+\langle
w,w'\rangle$.  It is the categorical product (and coproduct) in $\PreHilb$
and in $\Hilb$---a biproduct---and the zero space $0$ is a zero object.
On $\Hilb$ it is moreover a \emph{dagger biproduct}: each projection is
the dagger of the corresponding injection
\cite{heunen2019categories}.  The
canonical map $V^*\oplus W^*\to(V\oplus W)^*$, sending $(\xi,\eta)$ to
$(v,w)\mapsto\xi(v)+\eta(w)$, is a unitary isomorphism, and we use it
silently.  For a Hilbert space $V$, we write
\[
T^*V\coloneqq V\oplus V^*,
\]
call it the \emph{phase space} of $V$, and call $V^*$ the \emph{cotangent
space} $T^*_vV\coloneqq V^*$ at any point $v\in V$.  Combining the two
identifications above with reflexivity gives
$(T^*V)^*=(V\oplus V^*)^*\cong V^*\oplus V$, which we also use silently.

\paragraph{Sequence and function spaces.}
For a Hilbert space $B$ and a finite set $\Lambda$, we write
$B^{\oplus\Lambda}$ for the direct sum of copies of $B$ indexed by
$\Lambda$, and $\ell^2(\Lambda)\coloneqq\R^{\oplus\Lambda}$ with the
standard inner product.

For a measure space $(X,\mu)$ we write $L^2(X,\mu)$ for the Hilbert space
of \emph{square-integrable functions}: measurable functions
$f\colon X\to\R$ with $\int_X|f|^2\,d\mu<\infty$, taken modulo equality
$\mu$-almost everywhere, with inner product
$\langle f,g\rangle\coloneqq\int_Xfg\,d\mu$; completeness is the
Riesz--Fischer theorem \cite{brezis2011functional}.  For $Z\in\R_{>0}$ and
an interval $J\subseteq\R$, we write $Z\,dx$ for $Z$ times Lebesgue
measure on $J$; thus the \emph{weighted space} $L^2(J,Z\,dx)$ consists of
the square-integrable functions on $J$ with inner product
$\langle f,g\rangle=Z\int_Jfg\,dx$.  The case $J=[0,d)$ appears
throughout.

\paragraph{Dynamics and time change.}
For a monoid $(M,+,0)$, write $\Cat BM$ for the category with one object
whose endomorphism monoid is $M$.  Define
\[
\defineTerm{Dyn}_M\coloneqq[\Cat BM,\Hilb]:
\]
an object is a Hilbert space with an action of $M$ by bounded operators,
and a morphism is a bounded linear map commuting with the two actions.
Every monoid homomorphism $\mu\colon M\to M'$ induces a functor
$\Cat B\mu\colon\Cat BM\to\Cat BM'$, and restriction along it is a
functor
\[
\mu^*\coloneqq[\Cat B\mu,\Hilb]\colon\Dyn_{M'}\longrightarrow\Dyn_M,
\]
the \emph{time change} along $\mu$.  We write $\Dyn\coloneqq\Dyn_\N$,
where $\N=\{0,1,2,\ldots\}$: an object of $\Dyn$ is a pair $(X,F)$ of a
Hilbert space and a bounded endomorphism $F\colon X\to X$, its
\emph{transition}, and a morphism $\alpha\colon(X,F)\to(Y,G)$ is a
bounded linear map $\alpha\colon X\to Y$ satisfying
$\alpha\circ F=G\circ\alpha$.

\paragraph{Smooth maps.}
Let $V,W$ be Hilbert spaces and $f\colon V\to W$ a function.  The
\emph{Fr\'echet derivative} of $f$ at $v\in V$, when it exists, is the
bounded linear map $D_vf\colon V\to W$ satisfying
\[
\lim_{v'\to0}
\frac{\lVert f(v+v')-f(v)-D_vf(v')\rVert_W}{\lVert v'\rVert_V}=0.
\]
At most one linear map satisfies this: the difference $L$ of two
satisfying maps has $\lVert L(v')\rVert/\lVert v'\rVert\to0$ as $v'\to0$,
and by linearity this ratio is invariant under $v'\mapsto tv'$ for
$t>0$, so it vanishes for every $v'\neq0$, i.e.\ $L=0$.  The map $f$ is
\emph{smooth} if it has Fr\'echet derivatives of all orders, each depending
continuously on the basepoint \cite{cartan1971differential}.  Smooth maps
compose, and the chain rule $D_v(g\circ f)=D_{f(v)}g\circ D_vf$ holds.
Every bounded linear map $f$ is smooth with $D_vf=f$, as is every constant
map and every bounded affine map.  The pairing of smooth maps into a direct
sum is smooth.  For a smooth $U\colon V\to\R$ we write
$dU|_v\coloneqq D_vU\in V^*$, the derivative regarded as a bounded
functional, and call it the \emph{differential} of $U$ at $v$.  When
$V=V_1\oplus V_2$ is a direct sum, we write
$d_{V_1}U|_v\in V_1^*$ and $d_{V_2}U|_v\in V_2^*$ for the components of
$dU|_v$ under $(V_1\oplus V_2)^*\cong V_1^*\oplus V_2^*$, the
\emph{partial differentials}.

A \emph{continuous quadratic form} on a Hilbert space $V$ is a function
$U(v)\coloneqq\tfrac12 B(v,v)$ for a bounded symmetric bilinear
$B\colon V\times V\to\R$.  Every continuous quadratic form is smooth, with
differential the bounded linear map
\[
dU|_v=B(v,\blank)\in V^*,
\]
so $v\mapsto dU|_v$ is itself bounded linear.  Equivalently, by Riesz,
$U(v)=\tfrac12\langle A(v),v\rangle$ for the bounded operator
$A\coloneqq \riesz_V\circ dU\colon V\to V$, and symmetry of $B$ says
exactly that $A$ is \emph{self-adjoint}: $A^\dagger=A$.  We define
quadratic forms through bounded bilinear maps rather than---as one may in
finite dimensions \cite[Def.~5.3.3]{cdlm}%def.arrangement_terminology
---as linear functionals on a tensor square: on an infinite-dimensional
$V$, functionals on the Hilbert-space tensor square represent only the
\emph{Hilbert--Schmidt} forms \cite{conway1990functional}, and already
$U(v)=\tfrac12\lVert v\rVert^2$, whose operator is $A=\id$, is excluded.

A \emph{symmetric form} on a Hilbert space $A$ is a bounded linear map
$\lambda\colon A\to A^*$ with $\lambda^*=\lambda$ under reflexivity;
equivalently, by Riesz, a self-adjoint bounded operator
$\riesz_A\circ\lambda\colon A\to A$.  Symmetric forms on $A$ are exactly
the differentials of continuous quadratic forms:
$\lambda=d U|_{\blank}$ for $U(a)=\tfrac12\lambda(a)(a)$.

\section{The Hilbert arrangement datum}
\label{sec.hilbert_datum}

We now instantiate the framework.  Recall from
\cite[Def.~4.1.1]{cdlm}%def.rewiring_datum
that an adaptive-arrangement datum consists of a cartesian monoidal category
of spaces, a symmetric monoidal category of parameters, a strong monoidal
functor from parameters to spaces, and a monoidal potentials monad on
spaces.  The smooth datum of \cite[Def.~5.3.1]{cdlm}%def.potlens
takes finite-dimensional manifolds as spaces; here we take Hilbert spaces.

\subsection{Spaces}
\label{sec.hilbert_spaces}

\begin{definition}[Hilbert spaces with smooth maps]\label{def.hsm}
We define the category $\defineTerm{hsm}$ to have real Hilbert spaces as
objects and smooth maps (\cref{sec.background_hilbert}) as morphisms.
\end{definition}

\begin{proposition}\label{prop.hsm_cartesian}
$(\hsm,1,\oplus)$ is cartesian monoidal, where $1\coloneqq\R^0$.
\end{proposition}

\begin{proof}
The zero space $1=\R^0$ is terminal: the unique map $V\to1$ is constant,
hence smooth.  For Hilbert spaces $W_1,W_2$, the projections
$\pi_i\colon W_1\oplus W_2\to W_i$ are bounded linear, hence smooth.  Given
smooth $f_i\colon V\to W_i$, the pairing
$\langle f_1,f_2\rangle\colon V\to W_1\oplus W_2$ is smooth
(\cref{sec.background_hilbert}) and is the unique map whose composites with
the projections are the $f_i$.  Thus $\oplus$ is the categorical product,
and the monoidal structure it induces is cartesian.
\end{proof}

The pair $(\R,0,+)$ is a commutative monoid object in $\hsm$: addition
$\R\oplus\R\to\R$ is bounded linear and the unit $1\to\R$ is constant, so
both are smooth.  By \cite[Lem.~3.2.1]{cdlm},%lem.monoid_to_monad
it induces the monoidal \emph{writer monad} $\R\oplus\blank$ on $\hsm$,
which we abbreviate by its carrier $\R$.  This is our potentials monad: a
backward pass valued in $\R\oplus X$ is an ordinary backward pass valued in
$X$ together with a scalar potential.

\subsection{Parameters}
\label{sec.reactive_hilbert}

\begin{definition}[Reactive Hilbert space]\label{def.rhilb}
A \emph{reactive Hilbert space} is a pair $\rv V\coloneqq(V,\sharpR_V)$ of a
real Hilbert space $V$ and a bounded linear map
$\defineTerm{sharpR}_V\colon V^*\to V$, its \emph{reaction} or \emph{sharp}.  We call
$\rv V$ \emph{nondegenerate} if $\sharpR_V$ is an isomorphism, and
\emph{symmetric} if $\xi'(\sharpR_V\xi)=\xi(\sharpR_V\xi')$ for all
$\xi,\xi'\in V^*$.

We define the category $\defineTerm{rhilb}$ to have reactive Hilbert spaces
as objects and, as morphisms $\rv V\to\rv W$, the bounded linear
isomorphisms $\phi\colon V\To{\cong}W$ that are \emph{sharp-equivariant}:
\begin{equation}\label{eq.sharp_equivariance}
\phi\circ\sharpR_V\circ\phi^*=\sharpR_W
\end{equation}
as maps $W^*\to W$.
\end{definition}

The reactive vector spaces of \cite[Def.~5.2.1]{cdlm}%def.rvect
allow the reaction to vary smoothly with a basepoint $v\in V$; the harmonic
systems of this paper need only basepoint-independent reactions, so we
develop the constant case.  \Cref{eq.sharp_equivariance} is the constant
instance of the sharp-equivariance condition there.

\begin{proposition}\label{prop.rhilb_monoidal}
Direct sum makes $\rhilb$ symmetric monoidal, with unit $\rv 0\coloneqq(0,0)$
and
\[
\sharpR_{V\oplus W}\coloneqq\sharpR_V\oplus\sharpR_W
\colon V^*\oplus W^*\to V\oplus W
\]
under the identification $(V\oplus W)^*\cong V^*\oplus W^*$ of
\cref{sec.background_hilbert}.
\end{proposition}

\begin{proof}
A direct sum of bounded maps is bounded, so $\sharpR_{V\oplus W}$ is a
reaction.  If $\phi\colon\rv V\to\rv W$ and $\phi'\colon\rv V'\to\rv W'$ are
sharp-equivariant, then so is $\phi\oplus\phi'$, because transposes,
composites, and \cref{eq.sharp_equivariance} are all computed summandwise.
The coherence isomorphisms of $(\Hilb,0,\oplus)$ are sharp-equivariant for
the direct-sum reactions, again summandwise, so they lift to $\rhilb$.
\end{proof}

\begin{proposition}\label{prop.rhilb_forgetful}
The forgetful functor $\inc\colon\rhilb\to\hsm$, $(V,\sharpR_V)\mapsto V$, is
strong symmetric monoidal.
\end{proposition}

\begin{proof}
On both sides the tensor is $\oplus$ with unit the zero space
(\cref{prop.hsm_cartesian,prop.rhilb_monoidal}), and $\inc$ preserves both
on the nose; morphisms of $\rhilb$ are bounded linear, hence smooth.
\end{proof}

\subsection{The datum and its operad}
\label{sec.datum_operad}

\begin{definition}[Hilbert arrangement datum]\label{def.hilbert_datum}
The \emph{Hilbert arrangement datum} is
\[
\defineTerm{Hb}\coloneqq(\hsm,\rhilb,\inc,\R),
\]
an adaptive-arrangement datum in the sense of
\cite[Def.~4.1.1]{cdlm}:%def.rewiring_datum
$\hsm$ is cartesian monoidal (\cref{prop.hsm_cartesian}), $\rhilb$ is
symmetric monoidal (\cref{prop.rhilb_monoidal}), $\inc$ is strong symmetric
monoidal (\cref{prop.rhilb_forgetful}), and $\R$ abbreviates the writer
potentials monad $\R\oplus\blank$ (\cref{sec.hilbert_spaces}).  We write
$\defineTerm{harr}$ for its adaptive-arrangement operad.
\end{definition}

We unwind what \cite[Def.~4.1.1]{cdlm}%def.rewiring_datum
provides.  An object of $\harr$ is an \emph{interface} $\binom{A}{B}$ with
\emph{input space} $A:\hsm$ and \emph{output space} $B:\hsm$; the monoidal
unit is $I\coloneqq\binom{1}{1}$, and the tensor $\ltens$ direct-sums inputs
and outputs.  A morphism
$f\colon\bigl(\binom{A_1}{B_1},\ldots,\binom{A_K}{B_K}\bigr)\to\binom{A}{B}$,
which we call a \emph{Hilbert arrangement}, is a tuple
$f=(\rv Q,\outp f,\inpt f,U)$: a reactive Hilbert space
$\rv Q=(Q,\sharpR_Q)$, the \emph{parameter}, and smooth maps
\begin{equation}\label{eq.harr_morphism}
\begin{gathered}
\outp f\colon Q\oplus B_1\oplus\cdots\oplus B_K\to B,
\qquad
\inpt f\colon Q\oplus B_1\oplus\cdots\oplus B_K\oplus A
\to A_1\oplus\cdots\oplus A_K,\\
U\colon Q\oplus B_1\oplus\cdots\oplus B_K\oplus A\to\R,
\end{gathered}
\end{equation}
the \emph{output map}, the \emph{input map}, and the \emph{potential}.
(Products in $\hsm$ are direct sums, and in the writer instance the
backward pass splits into the input map and the potential.)  Composition
substitutes forward data into forward data and backward data into backward
data, tensors the parameters, and \emph{adds} the potentials of the
constituents after substitution; \cite[\S4.1]{cdlm}%rem.huge_wiring_diagram
draws the composite of two unary arrangements explicitly.  The identity
operation of $\binom AB$ is the static arrangement whose output and input
maps are identities.

\begin{definition}[Vocabulary]\label{def.vocabulary}
Following \cite[Def.~5.3.3]{cdlm},%def.arrangement_terminology
we say a Hilbert arrangement $f=(\rv Q,\outp f,\inpt f,U)$ is:
\emph{stateless} if $Q=\R^0$; \emph{potential-free} if $U=0$;
\emph{static} if both; a \emph{system} if its arity is $K=0$;
\emph{harmonic} if its reaction $\sharpR_Q$ is symmetric and its potential
$U$ is a continuous quadratic form (\cref{sec.background_hilbert}); and
\emph{linear} if its output map $\outp f$ and input map $\inpt f$ are
bounded linear.
\end{definition}

For a system, \cref{eq.harr_morphism} reduces to the parameter $\rv Q$,
an output map $\outp f\colon Q\to B$, and a potential
$U\colon Q\oplus A\to\R$; the input map is the unique map to $1$.

\section{Phase semantics}
\label{sec.hilbert_semantics}

Following the framework, semantics for $\harr$ takes two further
ingredients: a polynomial interpretation datum
\cite[Def.~4.2.1]{cdlm},%def.polynomial_interpretation
which casts spaces as polynomial interfaces and fixes what potentials do,
and an integrator \cite[Def.~4.3.4]{cdlm},%def.integrator
which turns the resulting effects into dynamics.  We build the Hilbert
instances of both and obtain the phase dynamics functor
$\Phiphase\colon\harr\to\pc$ (\cref{thm.phase_functor}).

\subsection{The cotangent functor}
\label{sec.cot_hilbert}

\begin{definition}[Cotangent functor]\label{def.cot_hilbert}
The \emph{cotangent functor} $\defineTerm{cot}\colon\hsm\to\poly$ sends a
Hilbert space $V$ to
\[
\cotof V\coloneqq\sum_{v\in V}\yon^{T^*_vV},
\]
where $T^*_vV=V^*$ (\cref{sec.background_hilbert}).  It sends a smooth map
$f\colon V\to W$ to the polynomial map $\cotof f\colon\cotof V\to\cotof W$
whose position map is $f$ and whose direction map at $v\in V$ is the
transpose of the Fr\'echet derivative,
$(D_vf)^*\colon T^*_{f(v)}W\to T^*_vV$.
\end{definition}

This is functorial: $D_v(\id_V)=\id_V$, whose transpose is the identity,
and for smooth $f\colon V\to W$, $g\colon W\to X$ the chain rule gives
$D_v(g\circ f)=D_{f(v)}g\circ D_vf$, so
$(D_v(g\circ f))^*=(D_vf)^*\circ(D_{f(v)}g)^*$, which is the direction map
of $\cotof g\circ\cotof f$.  This is the construction of
\cite[Def.~6.1.1]{cdlm}%def.cot
with manifolds replaced by Hilbert spaces, tangent maps by Fr\'echet
derivatives, and dual spaces by continuous duals; note that $\poly$ sees
only the underlying sets of the cotangent spaces.

\begin{proposition}\label{prop.cot_monoidal_hilbert}
The cotangent functor
$\cot\colon(\hsm,1,\oplus)\to(\poly,\yon,\otimes)$ is strong symmetric
monoidal.
\end{proposition}

\begin{proof}
The zero space has $1^*=0$, so $\cotof1=\sum_{\{0\}}\yon^{\{0\}}\cong\yon$.
For Hilbert spaces $V,W$, the identification
$(V\oplus W)^*\cong V^*\oplus W^*$ of \cref{sec.background_hilbert} gives
\[
\cotof{V\oplus W}
=\sum_{(v,w)\in V\oplus W}\yon^{(V\oplus W)^*}
\cong
\Bigl(\sum_{v\in V}\yon^{V^*}\Bigr)\otimes\Bigl(\sum_{w\in W}\yon^{W^*}\Bigr)
=\cotof V\otimes\cotof W
\]
by the definition of the Dirichlet product.  This isomorphism is natural:
for smooth $f\colon V\to V'$ and $g\colon W\to W'$, the product map
$f\oplus g$ has block-diagonal derivative
$D_{(v,w)}(f\oplus g)=D_vf\oplus D_wg$, whose transpose is
$(D_vf)^*\oplus(D_wg)^*$.  Coherence follows from the corresponding
coherence of $(1,\oplus)$ in $\hsm$, as in \cite[Prop.~6.1.2]{cdlm}.%prop.cot_monoidal
\end{proof}

\subsection{The Hilbert polynomial interpretation datum}
\label{sec.hilbert_interpretation}

Recall from \cite[Def.~4.2.1]{cdlm}%def.polynomial_interpretation
that a polynomial interpretation datum for a datum $D$ consists of an
interface functor (strong symmetric monoidal from spaces to $\poly$), a
semantic effect (a monoidal monad on $\poly$), an effect comparison
(a monoidal monad morphism carrying the syntactic effect to the semantic
one), and a potential algebra (a monoid with an algebra structure for the
semantic effect, which applies the effects).  For a writer potentials
monad, the canonical choices of semantic effect and comparison come with
the interface functor: the writer monad of the image of the carrier, and
the strong-monoidality comparison.

For the potential algebra we use the same object as the smooth instance.
The space $\R$ is finite-dimensional and $\R^*\cong\R$ by the standard
pairing, so the polynomial $\cotof\R=\sum_{a\in\R}\yon^{T^*_a\R}$ and its
commutative monoid structure (image of $(\R,0,+)$ under the strong monoidal
$\cot$) are literally those of \cite[\S6.2]{cdlm}.  Hence
\cite[Lem.~6.2.1]{cdlm}%lem.alpha_constant
applies unchanged: the $(\cotof\R\otimes\blank)$-monoid structures on
$\yon$ are exactly the constant covector fields $r\in\R$.  As there, we
take the canonical choice $r\coloneqq+1$ and write
$\defineTerm{potd}\colon\cotof\R\to\yon$ for the corresponding algebra map.  It is the
choice for which potentials contribute their differentials: for a smooth
map $U\colon V\to\R$ with $V:\hsm$, the composite
$\potd\circ\cotof U\colon\cotof V\to\yon$ is trivial on positions, and on
directions at $v\in V$ it selects
\begin{equation}\label{eq.potential_differential}
(D_vU)^*(+1)=dU|_v\in T^*_vV,
\end{equation}
so under the identification of maps $\cotof V\to\yon$ with covector fields
on $V$, it is the differential $dU$.

\begin{definition}[Hilbert polynomial interpretation datum]
\label{def.interp_hilbert}
The \emph{Hilbert polynomial interpretation datum} for $\Hb$ is
\[
\defineTerm{interph}\coloneqq
\bigl(\cot,\ \cotof\R\otimes\blank,\ \thetacomp,\ (\yon,\potd)\bigr),
\]
where $\cot$ is the cotangent functor
(\cref{def.cot_hilbert,prop.cot_monoidal_hilbert}), the semantic effect is
the writer monad of the commutative monoid $\cotof\R$, the comparison
$\defineTerm{thetacomp}\colon\cotof{\R\oplus\blank}\cong\cotof\R\otimes\cotof\blank$ is
induced by the strong monoidality of $\cot$, and the potential algebra is
$(\yon,\potd)$ as in \cref{eq.potential_differential}.
\end{definition}

\begin{corollary}\label{cor.interpretation_functor}
There is a normal lax symmetric monoidal functor
\[
\defineTerm{Phip}\colon\harr\longrightarrow\para{\cot}{\poly},
\]
acting on interfaces by
$\Phip_{\interph}\binom{A}{B}=\cotof B\otimes[\cotof A,\yon]$.
\end{corollary}

\begin{proof}
This is \cite[Thm.~4.2.2]{cdlm}%thm.poly_interpretation
applied to the datum $\Hb$ (\cref{def.hilbert_datum}) and the
interpretation datum $\interph$ (\cref{def.interp_hilbert}); it is normal
because the potential-algebra carrier is $\yon$.
\end{proof}

Here $\para{\cot}{\poly}$ denotes the Para construction for the action
$\rv Q\cdot p\coloneqq\cotof Q\otimes p$ of $\rhilb$ on $\poly$: an object
is a polynomial, and a morphism $p\to p'$ is a pair of a reactive Hilbert
space $\rv Q$ and a polynomial map $\cotof Q\otimes p\to p'$.

\subsection{The phase integrator}
\label{sec.phase_integrator}

Recall from \cite[Def.~4.3.4]{cdlm}%def.integrator
that for a strong symmetric monoidal $p\colon\cat Q\to\poly$, a
\emph{$p$-integrator} is a pair $(\Fn S,\fun{dyn})$ of a strong symmetric
monoidal state-space functor $\Fn S\colon\cat Q\to\smsetiso$ and a monoidal
natural transformation $\fun{dyn}\colon\Store\circ\Fn S\Rightarrow p$, and
that every $p$-integrator induces a normal lax symmetric monoidal functor
$\para{p}{\poly}\to\pc$ \cite[Prop.~4.3.5]{cdlm}.%prop.integrator_to_org
We build the phase integrator for $p\coloneqq\cot\circ\inc$, in three
steps: the phase-space endofunctor $T^*$, the Euler step $\chiE$, and the
readout $\nuro$.  Throughout we use the identifications
\[
\cotof V=\lvert V\rvert\,\yon^{V^*}
\qquad\text{and}\qquad
\cotof{T^*V}=\lvert V\oplus V^*\rvert\,\yon^{V^*\oplus V}
\]
of \cref{sec.background_hilbert}, where $\lvert\blank\rvert$ denotes the
underlying set.

\begin{proposition}\label{prop.TT_hilbert}
The assignments
\[
T^*\rv V\coloneqq\bigl(V\oplus V^*,\ \sharpS_{T^*V}\bigr),
\qquad
T^*\phi\coloneqq\phi\oplus(\phi^*)^{-1},
\]
where the \emph{canonical symplectic sharp} is the bounded map
\begin{equation}\label{eq.symplectic_sharp}
\defineTerm{sharpS}_{T^*V}\colon V^*\oplus V\to V\oplus V^*,
\qquad
(\xi,v)\mapsto(v,-\xi),
\end{equation}
define a strong symmetric monoidal endofunctor $T^*\colon\rhilb\to\rhilb$.
\end{proposition}

\begin{proof}
We first check that $T^*\phi$ is a morphism of $\rhilb$ for every
$\phi\colon\rv V\to\rv W$.  It is a bounded linear isomorphism because
$\phi$ and $(\phi^*)^{-1}$ are.  For sharp-equivariance
\cref{eq.sharp_equivariance} we must show
$T^*\phi\circ\sharpS_{T^*V}\circ(T^*\phi)^*
=\sharpS_{T^*W}$.  Under the identifications of
\cref{sec.background_hilbert}, the transpose of
$T^*\phi=\phi\oplus(\phi^*)^{-1}$ is
$(T^*\phi)^*=\phi^*\oplus\phi^{-1}$ as a map
$W^*\oplus W\to V^*\oplus V$, using reflexivity and
$((\phi^*)^{-1})^*=(\phi^{**})^{-1}=\phi^{-1}$.  So for
$(\eta,w)\in W^*\oplus W$:
\[
(\eta,w)
\;\xmapsto{(T^*\phi)^*}\;
(\phi^*\eta,\ \phi^{-1}w)
\;\xmapsto{\sharpS_{T^*V}}\;
(\phi^{-1}w,\ -\phi^*\eta)
\;\xmapsto{T^*\phi}\;
\bigl(w,\ -\eta\bigr)
=\sharpS_{T^*W}(\eta,w).
\]
Note that the sharp-equivariance of $\phi$ itself was not needed: the
symplectic sharp is the same formula for every reactive Hilbert space, as
in \cite[Prop.~6.3.4]{cdlm}.%prop.TT_endofunctors
Preservation of identities and composites is inherited from $\Hilb$ and
from $(g\circ f)^*=f^*\circ g^*$.

For strong monoidality, the unitor is the identity, since
$T^*\rv 0=0$, and the productor is the coordinate swap
\[
(T^*)_2\colon T^*V\oplus T^*W\To{\cong}T^*(V\oplus W),
\qquad
\bigl((v,\xi),(w,\eta)\bigr)\mapsto\bigl((v,w),(\xi,\eta)\bigr).
\]
By \cref{eq.symplectic_sharp}, the symplectic sharp of $T^*(V\oplus W)$
sends $(\xi,\eta,v,w)\mapsto(v,w,-\xi,-\eta)$; conjugating by $(T^*)_2$
gives $(\xi,v,\eta,w)\mapsto(v,-\xi,w,-\eta)$, which is
$\sharpS_{T^*V}\oplus\sharpS_{T^*W}$, the direct-sum
sharp of \cref{prop.rhilb_monoidal}.  Symmetry and coherence are inherited
from $(\Hilb,0,\oplus)$, as in \cite[Prop.~6.3.5]{cdlm}.%prop.TT_monoidal
\end{proof}

Since every morphism of $\rhilb$ is an isomorphism, the underlying-set
functors $\lvert\blank\rvert$ and $\lvert T^*\blank\rvert$ land in
$\smsetiso$, and both are strong symmetric monoidal:
$\lvert V\oplus W\rvert\cong\lvert V\rvert\times\lvert W\rvert$.

\begin{proposition}[Euler step]\label{prop.euler_step_hilbert}
There is a monoidal natural transformation
$\defineTerm{chiE}\colon\Store\circ\lvert\blank\rvert\Rightarrow\cot\circ\inc$ of
functors $\rhilb\to\poly$ whose component
$\chiE_{\rv V}\colon\Store(\lvert V\rvert)\to\cotof V$ is the identity on
positions and, on directions at $v\in V$, is
\[
T^*_vV=V^*\to\lvert V\rvert,
\qquad
\xi\mapsto v+\sharpR_V(\xi).
\]
When $\rv V$ is nondegenerate, $\chiE_{\rv V}$ is an isomorphism.
\end{proposition}

\begin{proof}
Both composites send $\rv V$ to a polynomial with position set
$\lvert V\rvert$, so the identity is a valid position map, and the stated
direction maps assemble into a polynomial map
$\Store(\lvert V\rvert)=\lvert V\rvert\,\yon^{\lvert V\rvert}\to
\lvert V\rvert\,\yon^{V^*}$.  For naturality at
$\phi\colon\rv V\to\rv W$, the position squares commute by definition; on
directions at $v\in V$, the two composites
$T^*_{\phi(v)}W\to\lvert V\rvert$ are
\[
\xi\;\xmapsto{\phi^*}\;\phi^*\xi\;\xmapsto{}\;v+\sharpR_V(\phi^*\xi)
\qquad\text{and}\qquad
\xi\;\xmapsto{}\;\phi(v)+\sharpR_W(\xi)\;\xmapsto{\phi^{-1}}\;
v+\phi^{-1}(\sharpR_W\xi),
\]
along $\chiE_{\rv V}\circ\cotof\phi$ and
$\Store(\phi)\circ\chiE_{\rv W}$ respectively.  These agree if and only if
$\sharpR_V\circ\phi^*=\phi^{-1}\circ\sharpR_W$, which is sharp-equivariance
\cref{eq.sharp_equivariance}.  Monoidality holds because
$\sharpR_{V\oplus W}=\sharpR_V\oplus\sharpR_W$
(\cref{prop.rhilb_monoidal}) and translation acts coordinatewise.  When
$\sharpR_V$ is an isomorphism, so is each direction map, hence so is
$\chiE_{\rv V}$.  This is \cite[Prop.~6.3.2]{cdlm}%prop.rvect_polynomial
with constant reactions and bounded sharps.
\end{proof}

\begin{proposition}[Phase readout]\label{prop.nu_hilbert}
There is a monoidal natural transformation
$\defineTerm{nuro}\colon\cot\circ\inc\circ\,T^*\Rightarrow\cot\circ\inc$ of functors
$\rhilb\to\poly$ whose component
$\nuro_{\rv V}\colon\cotof{T^*V}\to\cotof V$ has position map the
\emph{presented position}
\begin{equation}\label{eq.presented_position_abstract}
(v,\xi)\mapsto\tilde v\coloneqq v+\sharpR_V(\xi)
\end{equation}
and direction map at $(v,\xi)$ given by
\[
V^*\to V^*\oplus V,
\qquad
\xi'\mapsto\bigl(\xi',\ \sharpR_V(\xi)\bigr).
\]
\end{proposition}

\begin{proof}
The position map is bounded linear, and the direction maps land in
$V^*\oplus V\cong T^*_{(v,\xi)}(T^*V)$, so the components are polynomial
maps.  For naturality at $\phi\colon\rv V\to\rv W$ we compare
$\cotof\phi\circ\nuro_{\rv V}$ with $\nuro_{\rv W}\circ\cotof{T^*\phi}$.  On
positions, using sharp-equivariance \cref{eq.sharp_equivariance} in the
middle step,
\[
\phi\bigl(v+\sharpR_V\xi\bigr)
=\phi(v)+\phi\bigl(\sharpR_V\xi\bigr)
=\phi(v)+\sharpR_W\bigl((\phi^*)^{-1}\xi\bigr),
\]
which is the presented position of
$T^*\phi(v,\xi)=(\phi v,(\phi^*)^{-1}\xi)$.  On directions at $(v,\xi)$,
take $\xi'\in W^*$.  Along $\cotof\phi\circ\nuro_{\rv V}$ it maps to
$\bigl(\phi^*\xi',\ \sharpR_V\xi\bigr)$.  Along
$\nuro_{\rv W}\circ\cotof{T^*\phi}$ it maps first to
$\bigl(\xi',\sharpR_W(\phi^*)^{-1}\xi\bigr)$ and then, along the transpose
$(T^*\phi)^*=\phi^*\oplus\phi^{-1}$ (proof of \cref{prop.TT_hilbert}), to
$\bigl(\phi^*\xi',\ \phi^{-1}\sharpR_W(\phi^*)^{-1}\xi\bigr)$; and
$\phi^{-1}\circ\sharpR_W\circ(\phi^*)^{-1}=\sharpR_V$ by
\cref{eq.sharp_equivariance}.  Monoidality holds coordinatewise, as in
\cref{prop.euler_step_hilbert}.  This transformation is the composite
readout $\nuro_{\exp,\beta}$ of \cite[Lem.~6.3.16]{cdlm},%lem.combined_update
assembled from the exponential readout \cite[Lem.~6.3.9]{cdlm}%lem.readout_examples
and the kinetic $1$-form \cite[Ex.~6.3.13]{cdlm};%ex.kinetic_one_form
we have verified it directly.
\end{proof}

\begin{definition}[Phase integrator]\label{def.phase_integrator_hilbert}
The \emph{phase dynamics} is the monoidal natural transformation
\[
\defineTerm{dynphase}\colon
\Store\circ\lvert T^*\blank\rvert
\xRightarrow{\ \chiE_{T^*\blank}\ }
\cot\circ\inc\circ\,T^*
\xRightarrow{\ \nuro\ }
\cot\circ\inc,
\]
the Euler step on phase space (\cref{prop.euler_step_hilbert} whiskered by
the strong monoidal $T^*$ of \cref{prop.TT_hilbert}) followed by the phase
readout (\cref{prop.nu_hilbert}).  The \emph{phase integrator} is the
$(\cot\circ\inc)$-integrator
\[
\phase\coloneqq\bigl(\lvert T^*\blank\rvert,\ \dynphase\bigr)
\]
in the sense of \cite[Def.~4.3.4]{cdlm}.%def.integrator
\end{definition}

We record the resulting update.  At $\rv V$, the direction map of
$\chiE_{T^*\rv V}$ adds the symplectic sharp of the covector assembled by
$\nuro$: by \cref{eq.symplectic_sharp},
$\sharpS_{T^*V}(\xi',\sharpR_V\xi)=(\sharpR_V\xi,\,-\xi')$, so on
directions
\begin{equation}\label{eq.phase_update}
\lvert T^*V\rvert\times V^*\to\lvert T^*V\rvert,
\qquad
\bigl((v,\xi),\xi'\bigr)\mapsto
\bigl(v+\sharpR_V(\xi),\ \xi-\xi'\bigr),
\end{equation}
while the position reported is the presented position
\cref{eq.presented_position_abstract}.

\begin{theorem}[Phase dynamics functor]\label{thm.phase_functor}
There is a normal operad functor
\[
\defineTerm{Phiphase}\coloneqq\Phi_{\interph,\phase}
\colon\harr\longrightarrow\pc,
\]
with
$\Phiphase\binom{A}{B}=\cotof B\otimes[\cotof A,\yon]$ on interfaces and
$\Phiphase(I)\cong\yon$.
\end{theorem}

\begin{proof}
This is \cite[Thm.~4.3.6]{cdlm}%thm.dynamics_functor
applied to the interpretation datum $\interph$
(\cref{def.interp_hilbert}) and the integrator $\phase$
(\cref{def.phase_integrator_hilbert}); normality, hence
$\Phiphase(I)\cong\yon$, holds because the potential-algebra carrier is
$\yon$.
\end{proof}

The physical constants of this paper enter through the parameter and the
potential rather than through the functor: for an \emph{impedance}
$Z\in\R_{>0}$ and a \emph{cell duration} $h\in\R_{>0}$, a unit mass of
size $Zh$ has reaction $\riesz/(Zh)$, a unit bond of stiffness $Z/h$ has
potential $\frac{Z}{2h}(a-q)^2$, and executing one $\pc$-tick per duration
$h$ multiplies both rates by $h$.  The systems of
\cref{sec.cell} are these tuned combinations: reaction $\riesz/Z$ and
stiffness $Z$, with no residual $h$.  The duration $h$ reappears only in
the semantics of packets (\cref{sec.queues}), where it weights the packet
inner products and calibrates subdivision.

\section{Open systems and ported machines}
\label{sec.open}

The interfaces of $\harr$ carry inputs as well as outputs, and this
section equips the semantics to use them.  On the syntactic side, open
harmonic-linear systems form a semicategory $\Open$: serial gluing is
substitution into a seam wiring, and the glued potential is the sum of
the constituents' potentials across the seam.  On the semantic side,
bounded linear machines with two-sided ports form a semicategory
$\Lin$.  The refinement $\Phiopen\colon\Open\to\Lin$ is the unique
semifunctor compatible with $\Phiphase$.  Machines also compare across
ports: \cref{sec.machine_colimits} equips them with morphisms over
port embeddings and with tower colimits.

Throughout, composites of machines carry direct sums of state spaces,
and we suppress the associators of $\oplus$, as one suppresses the
product associators of carriers in $\pc$.

\subsection{Ports and open systems}
\label{sec.ports}

\begin{definition}[Port]\label{def.port}
A \emph{port} is a pair $(A,\lambda)$ of a real Hilbert space $A$, its
\emph{carrier}, and a symmetric form $\lambda\colon A\to A^*$
(\cref{sec.background_hilbert}), its \emph{impedance form}.
\end{definition}

A port is the boundary datum of an open harmonic system: the carrier
types the values that cross the boundary, and the impedance form is the
linear part of the covector field that the system returns along it.

\begin{definition}[Seam]\label{def.seam}
For Hilbert spaces $A,B,C$, the \emph{seam} is the static arrangement
\[
\defineTerm{seam}_{A,B,C}
\colon
\Bigl(\binom AB,\ \binom BC\Bigr)
\longrightarrow
\binom AC
\]
whose output map is $(b,c)\mapsto c$ and whose input map is
$(b,c,a)\mapsto(a,b)$.
\end{definition}

In the monoidal reading of the operad
(\cref{sec.background_categories}), substituting a pair of systems
into the seam tensors them and composes with the seam morphism
$\binom AB\ltens\binom BC\to\binom AC$.  The routing is serial: the
outer input feeds the first box, the first box's output feeds the
second, and the second box's output is the outer output; the backward
routing is the substitution discipline of $\harr$, not further data.

\begin{lemma}[Seams associate]\label{lem.seams_associate}
For Hilbert spaces $A,B,C,D$, the two substitutions of one seam into
another agree:
\[
\seam_{A,C,D}\bigl(\seam_{A,B,C},\ \id\bigr)
=
\seam_{A,B,D}\bigl(\id,\ \seam_{B,C,D}\bigr)
\colon
\Bigl(\binom AB,\ \binom BC,\ \binom CD\Bigr)
\longrightarrow
\binom AD,
\]
both being the static arrangement with output map $(b,c,d)\mapsto d$
and input map $(b,c,d,a)\mapsto(a,b,c)$.
\end{lemma}

\begin{proof}
A substitution of static arrangements is static with the composed
routing (\cref{sec.datum_operad}): parameters tensor to $\rv0$,
potentials add to $0$, and forward and input maps substitute.  On the
left: the outer seam's output is its second box's output $d$, and its
input map feeds $a$ to the inner seam and $c$ to the second box; the
inner seam feeds $a$ to the first box and $b$ to its own second box.
On the right: the output is the inner seam's output $d$; the outer
seam feeds $a$ to the first box and $b$ to the inner seam, which feeds
$b$ to its first box and $c$ to its second.  Both routings are
$(b,c,d,a)\mapsto(a,b,c)$.
\end{proof}

\begin{definition}[Open systems]\label{def.open}
The semicategory $\defineTerm{Open}$ has as objects the ports.  A
morphism $f\colon(A,\lambda)\to(B,\mu)$ is a harmonic linear system
(\cref{def.vocabulary}) $f=(\rv Q,\outp f,U)\colon I\to\binom{A}{B}$
in $\harr$ whose potential satisfies the \emph{impedance condition}
\begin{equation}\label{eq.impedance_condition}
d_AU|_{(0,a)}=\lambda(a)
\qquad(a\in A).
\end{equation}
Composition is substitution into the seam,
\[
g\circ f\coloneqq\seam_{A,B,C}(f,g),
\]
and the direct sum of ports,
$(A,\lambda)\oplus(A',\lambda')\coloneqq(A\oplus A',\lambda\oplus\lambda')$,
and of systems makes $\Open$ symmetric semigroupal.
\end{definition}

Associativity is inherited from the operad: by the associativity of
substitution,
$h\circ(g\circ f)$ and $(h\circ g)\circ f$ are the substitutions of
$(f,g,h)$ into the two sides of \cref{lem.seams_associate}.  That the
composite is again a morphism of $\Open$---harmonic, linear, and
satisfying the impedance condition---is part of the next proposition,
which also records the composite in components; the middle port's form
$\mu$ places no condition on the composite and is instead consumed by
the semantics (\cref{def.ported_machine}).

\begin{proposition}[Composite data]\label{prop.open_composite}
For composable $f\colon(A,\lambda)\to(B,\mu)$ and
$g\colon(B,\mu)\to(C,\nu)$ in $\Open$,
\begin{equation}\label{eq.open_composition}
g\circ f
=
\Bigl(
\rv Q_f\oplus\rv Q_g,\
\outp g\circ\pi_{Q_g},\
U_{g\circ f}
\Bigr),
\qquad
U_{g\circ f}\bigl((q_f,q_g),a\bigr)
=
U_f(q_f,a)+U_g\bigl(q_g,\outp f(q_f)\bigr),
\end{equation}
and $g\circ f$ is a morphism $(A,\lambda)\to(C,\nu)$ of $\Open$.
\end{proposition}

\begin{proof}
We unwind the substitution described in \cref{sec.datum_operad}.  The
substituted arrangement has parameter the tensor
$\rv Q_f\oplus\rv Q_g$ of the constituents' parameters, the seam being
stateless.  Its output map substitutes the inner outputs into the
seam's output map: $(q_f,q_g)\mapsto\outp g(q_g)$.  Its potential adds
the constituents' potentials after substituting forward data: the seam
contributes zero; $f$ receives the outer input $a$, contributing
$U_f(q_f,a)$; and $g$ receives $f$'s output, contributing
$U_g(q_g,\outp f(q_f))$.  This is \cref{eq.open_composition}.

The composite is a morphism of $\Open$: the parameter
$\rv Q_f\oplus\rv Q_g$ has symmetric reaction
(\cref{prop.rhilb_monoidal}); the output map is bounded linear; and
$U_{g\circ f}$ is a continuous quadratic form, being a sum of
continuous quadratic forms precomposed with bounded linear maps.  For
the impedance condition, the second summand of $U_{g\circ f}$ does not
depend on $a$, so
$d_AU_{g\circ f}|_{(0,a)}=d_AU_f|_{(0,a)}=\lambda(a)$.
\end{proof}

There are no identity morphisms when $A\neq0$: an identity would
require an output map whose value is the incoming input, and output
maps of systems read only the parameter \cref{eq.harr_morphism}.  The
same absence recurs in the semantics, and \cref{sec.divisible} realizes
it as the nonexistence of strings of duration zero.

\subsection{Ported machines}
\label{sec.ported_machines}

The interface polynomial $\Phiphase\binom AB=\cotof B\otimes[\cotof
A,\yon]$ (\cref{thm.phase_functor}) has more positions than outputs: it
also presents a covector field on the input space.

\begin{lemma}[Field interface]\label{lem.field_interface}
For $A:\Hilb$, write $\defineTerm{Fl}(A)$ for the set of functions
$\lvert A\rvert\to A^*$, the \emph{covector fields} on $A$.  Then
\[
[\cotof A,\yon]
\;\cong\;
\Fl(A)\,\yon^{\lvert A\rvert}:
\]
a position is a covector field on $A$, and a direction at every
position is a point of $A$.
\end{lemma}

\begin{proof}
For every polynomial $q$, the hom-tensor adjunction of $\poly$ gives
$\poly\bigl(q,[\cotof A,\yon]\bigr)\cong\poly(q\otimes\cotof A,\yon)$.
A map $q\otimes\cotof A\to\yon$ is trivial on positions and selects, at
each position $(j,a)$ of $q\otimes\cotof A$, a direction of
$q\otimes\cotof A$ there, i.e.\ an element of $q[j]\times T^*_aA$;
equivalently, for each $j$, a pair of a function
$\lvert A\rvert\to A^*$ and a function $\lvert A\rvert\to q[j]$.  This
is exactly a map $q\to\Fl(A)\,\yon^{\lvert A\rvert}$: a position
map $j\mapsto\kappa_j$ into the fields together with direction maps
$\lvert A\rvert\to q[j]$.  The bijection is natural in $q$, so the
Yoneda lemma gives the isomorphism.
\end{proof}

\begin{definition}[Affine embedding]\label{def.affine_embedding}
For a port $(A,\lambda)$, the \emph{affine embedding} is the polynomial
map
\[
\iota_{(A,\lambda)}
\colon
\lvert A^*\rvert\,\yon^{\lvert A\rvert}
\longrightarrow
\Fl(A)\,\yon^{\lvert A\rvert}
\cong[\cotof A,\yon],
\]
sending the position $\zeta\in A^*$ to the affine field
$a\mapsto\lambda(a)+\zeta$, with identity direction maps.  It is
injective on positions and bijective on directions.
\end{definition}

\begin{definition}[Ported machines]\label{def.ported_machine}
The semicategory $\defineTerm{Lin}$ has as objects the ports.  Its
hom-category $\Lin\bigl((A,\lambda),(B,\mu)\bigr)$ has as objects the
\emph{ported machines} $(S,f^+,f^-)$: a Hilbert space $S$, the
\emph{state space}, with bounded linear maps
\[
f^+\colon S\to B\oplus A^*,
\qquad
f^-\colon S\oplus A\oplus B^*\to S,
\]
the \emph{readout} and the \emph{update}; a morphism
$\alpha\colon(S,f^+,f^-)\to(\bar S,\bar f^+,\bar f^-)$ is a bounded
linear map $\alpha\colon S\to\bar S$ with
\[
\bar f^+\circ\alpha=f^+
\qquad\text{and}\qquad
\bar f^-\circ(\alpha\oplus\id_A\oplus\id_{B^*})=\alpha\circ f^-.
\]
We write $f^+=\langle f^+_B,f^+_{A^*}\rangle$ for the components of
the readout, the \emph{value} and the \emph{offset}.  The composite of
$M=(S_M,f^+,f^-)\colon(A,\lambda)\to(B,\mu)$ and
$N=(S_N,g^+,g^-)\colon(B,\mu)\to(C,\nu)$ is the ported machine with
state space $S_M\oplus S_N$ and
\begin{equation}\label{eq.machine_composition}
\begin{gathered}
(N\circ M)^+(s,s')
\coloneqq
\bigl(g^+_C(s'),\ f^+_{A^*}(s)\bigr),\\
(N\circ M)^-\bigl((s,s'),a,\theta\bigr)
\coloneqq
\Bigl(
f^-\bigl(s,\ a,\ \mu(f^+_B s)+g^+_{B^*}(s')\bigr),\
g^-\bigl(s',\ f^+_B(s),\ \theta\bigr)
\Bigr);
\end{gathered}
\end{equation}
on morphisms, $\alpha\circ\beta\coloneqq\alpha\oplus\beta$.
\end{definition}

The composition formula is the seam in machine form: $N$ receives
$M$'s value as its input, and $M$ receives, as its incoming covector,
the field that $N$ returns along its own input port---the impedance
part $\mu(f^+_Bs)$ contributed by the typing, the offset part
$g^+_{B^*}(s')$ contributed by $N$'s state.  All maps in
\cref{eq.machine_composition} are composites and sums of bounded
linear maps, morphisms compose componentwise, and associativity is a
direct computation using the associativity of $\oplus$: in a triple
composite, each seam inserts its own middle impedance form, and the
three state spaces update side by side.  Direct sum of ports and of
machines makes $\Lin$ symmetric semigroupal.  As in $\Open$, there are
no identity 1-cells when $A\neq0$: a readout depends only on the
state, so no machine can emit its current input.

\begin{definition}[Underlying 1-cells]\label{def.ufun}
For a ported machine $M=(S,f^+,f^-)\colon(A,\lambda)\to(B,\mu)$, define
$\defineTerm{Ufun}(M)\in\pc\bigl(\yon,\ \cotof B\otimes[\cotof A,\yon]\bigr)$
to be the $1$-cell with carrier $\lvert S\rvert$ whose position at
$s\in S$ is
\[
\Bigl(f^+_B(s),\
a\mapsto\lambda(a)+f^+_{A^*}(s)\Bigr)
\]
under \cref{lem.field_interface}, and whose direction map at $s$ sends
$(\eta,a)\in B^*\times\lvert A\rvert$ to the next state
$f^-(s,a,\eta)$.
\end{definition}

The assignment $\Ufun$ is faithful and injective on machines: the
value, the offset (recovered from the affine field by evaluating at
$0$), and the update are all visible in the $1$-cell, and the
commutation conditions on $2$-cells match the coalgebra-morphism
conditions.  Its image consists of the $1$-cells that factor through
the affine embeddings $\iota_{(A,\lambda)}$ with bounded linear
structure maps.

\subsection{The open linear semantics}
\label{sec.open_semantics}

We now compute $\Phiphase$ on systems and on the seam, and assemble
the results into the semantics $\Phiopen$.

\begin{lemma}[System unwinding]\label{lem.system_unwinding}
Let $f=(\rv Q,\outp f,U)\colon I\to\binom AB$ be a system in $\harr$.
Then $\Phiphase(f)\in\pc\bigl(\yon,\cotof B\otimes[\cotof A,\yon]\bigr)$
is the $1$-cell with carrier $\lvert T^*Q\rvert$ whose position at
$(q,\xi)$ is
\[
\Bigl(\outp f(\tilde q),\
a\mapsto d_AU|_{(\tilde q,a)}\Bigr),
\qquad
\tilde q\coloneqq q+\sharpR_Q(\xi)
\]
the presented position \cref{eq.presented_position_abstract}, and whose
direction map at $(q,\xi)$ sends $(\eta,a)\in B^*\times\lvert A\rvert$
to the next state
\[
\Bigl(\tilde q,\
\xi-(D_{\tilde q}\outp f)^*\eta-d_QU|_{(\tilde q,a)}\Bigr).
\]
\end{lemma}

\begin{proof}
This is the unwinding of the phase coalgebra
\cite[\S6.6.1]{cdlm},%sec.phase_coalgebra
whose component formulas
\cite[Eqs.~(53), (56), (62)]{cdlm}%eqn.outpn, eqn.covector_triple, eqn.state_update
depend only on the data $\chiE$, $\nuro$, $\potd$ shared verbatim by
the Hilbert instance
(\cref{prop.euler_step_hilbert,prop.nu_hilbert,eq.potential_differential}).
At a state $(q,\xi)$, the integrator presents the parameter position
$\tilde q$ \cref{eq.presented_position_abstract}; the forward data,
evaluated there, gives the output value $\outp f(\tilde q)$ and, since
the arity is zero, no inner routing.  The backward pass at received
input $a$ assembles the total covector
\[
\omega
\coloneqq
(D_{\tilde q}\outp f)^*\eta+dU|_{(\tilde q,a)}
\ \in\ Q^*\oplus A^*,
\]
the pullback sum of \cite[Eq.~(56)]{cdlm}:%eqn.covector_triple
its $A^*$-component $d_AU|_{(\tilde q,a)}$ is returned along the input
interface---as a function of the yet-to-arrive $a$, it is the field
component of the position, by \cite[Eq.~(53)]{cdlm}%eqn.outpn
under \cref{lem.field_interface}---and its $Q^*$-component enters the
integrator update \cref{eq.phase_update}, giving the stated next
state by \cite[Eq.~(62)]{cdlm}.%eqn.state_update
\end{proof}

\begin{lemma}[Seam unwinding]\label{lem.seam_unwinding}
$\Phiphase(\seam_{A,B,C})$ is the stateless $1$-cell
(\cref{lem.stateless}) of the polynomial map
\[
\Bigl(\cotof B\otimes[\cotof A,\yon]\Bigr)
\otimes
\Bigl(\cotof C\otimes[\cotof B,\yon]\Bigr)
\longrightarrow
\cotof C\otimes[\cotof A,\yon]
\]
whose position map sends $\bigl((b,\kappa),(c,\kappa')\bigr)$ to
$(c,\kappa)$ and whose direction map there sends
$(\theta,a)\in C^*\times\lvert A\rvert$ to
$\bigl((\kappa'(b),a),\ (\theta,b)\bigr)$.
\end{lemma}

\begin{proof}
The seam is static, so its parameter is $\rv0$, the integrator carrier
is a single state, and the potential contributes nothing: the same
unwinding as in \cref{lem.system_unwinding}, now with inner interfaces,
reduces to the forward and backward routing of the wiring
\cite[Eqs.~(53), (56), (62)]{cdlm}.%eqn.outpn, eqn.covector_triple, eqn.state_update
Forward: the outer output is the seam's output map applied to the
inner values, $(b,c)\mapsto c$, and the outer field is the pullback of
the inner fields along the input map's dependence on $a$: the input
map $(b,c,a)\mapsto(a,b)$ feeds $a$ to the first box, so the outer
field is $\kappa$.  Backward, at outer direction $(\theta,a)$: the
input map routes $a$ to the first box and $b$ to the second; the
covector returned by the second box along its input, $\kappa'$
evaluated at its received input $b$, pulls back along the input map to
the first box's output interface; and the outer covector $\theta$
pulls back along the output map to the second box's output interface.
\end{proof}

\begin{definition}[Machine of a system]\label{def.machine_of_system}
Let $f=(\rv Q,\outp f,U)\colon(A,\lambda)\to(B,\mu)$ be a morphism of
$\Open$.  Its \emph{machine} $M_f$ is the ported machine with state
space $T^*Q=Q\oplus Q^*$ and, writing
$\tilde q\coloneqq q+\sharpR_Q(\xi)$,
\begin{equation}\label{eq.machine_of_system}
\begin{gathered}
M_f^+(q,\xi)
\coloneqq
\bigl(\outp f(\tilde q),\ d_AU|_{(\tilde q,0)}\bigr),\\
M_f^-\bigl((q,\xi),a,\eta\bigr)
\coloneqq
\bigl(\tilde q,\ \xi-(\outp f)^*\eta-d_QU|_{(\tilde q,a)}\bigr).
\end{gathered}
\end{equation}
\end{definition}

Every map in \cref{eq.machine_of_system} is a composite of the bounded
linear maps $\outp f$, $(\outp f)^*$, $\sharpR_Q$, and
$dU|_{\blank}$---bounded linear in its basepoint because $U$ is a
continuous quadratic form (\cref{sec.background_hilbert})---so $M_f$
is a ported machine.

\begin{theorem}[Open linear semantics]\label{thm.open_semantics}
There is a unique semifunctor
\[
\defineTerm{Phiopen}\colon\Open\longrightarrow\Lin,
\]
the identity on ports, satisfying
$\Ufun\circ\Phiopen=\Phiphase$ on morphisms.  Its value on $f$ is the
machine $M_f$ of \cref{def.machine_of_system}.
\end{theorem}

\begin{proof}
We first show $\Ufun(M_f)=\Phiphase(f)$ for every
$f\colon(A,\lambda)\to(B,\mu)$.  Both $1$-cells have carrier
$\lvert T^*Q\rvert$.  Since $U$ is a continuous quadratic form, its
differential is linear in the basepoint, so the field presented by
\cref{lem.system_unwinding} at $(q,\xi)$ splits as
\[
d_AU|_{(\tilde q,a)}
=
d_AU|_{(\tilde q,0)}+d_AU|_{(0,a)}
=
d_AU|_{(\tilde q,0)}+\lambda(a),
\]
using the impedance condition \cref{eq.impedance_condition}: it is the
affine field of the offset $d_AU|_{(\tilde q,0)}=M_f^+$'s offset
component, exactly the position of $\Ufun(M_f)$
(\cref{def.ufun}).  The values agree, $\outp f$ being linear; and the
direction maps agree because $D_{\tilde q}\outp f=\outp f$.

Now let $f\colon(A,\lambda)\to(B,\mu)$ and $g\colon(B,\mu)\to(C,\nu)$
be composable in $\Open$.  Composition in $\Open$ is substitution into
the seam (\cref{def.open}), so operad functoriality
(\cref{thm.phase_functor}) gives
\[
\Phiphase(g\circ f)
=\Phiphase\bigl(\seam_{A,B,C}(f,g)\bigr)
=\Phiphase(\seam_{A,B,C})
\circ\bigl(\Phiphase(f)\otimes\Phiphase(g)\bigr).
\]
Substituting $\Phiphase(f)=\Ufun(M_f)$ and
$\Phiphase(g)=\Ufun(M_g)$ and unwinding the right side by
\cref{lem.seam_unwinding}: the composite $1$-cell has carrier
$\lvert T^*Q_f\rvert\times\lvert T^*Q_g\rvert$; its position at
$(s,s')$ is $\bigl(g$'s value$,\ f$'s field$\bigr)$; and its direction
map at $(\theta,a)$ feeds $a$ and the evaluated field
$\kappa'(f^+_Bs)=\mu(f^+_Bs)+g^+_{B^*}(s')$ to $M_f$, and $f$'s value
and $\theta$ to $M_g$.  This is exactly
$\Ufun(M_g\circ M_f)$ by
\cref{eq.machine_composition,def.ufun}.  Applying the first paragraph
to $g\circ f$ gives
$\Ufun(M_{g\circ f})=\Phiphase(g\circ f)=\Ufun(M_g\circ M_f)$, and
$\Ufun$ is injective on machines (\cref{def.ufun}), so
$M_{g\circ f}=M_g\circ M_f$: the assignment $f\mapsto M_f$ is a
semifunctor.  Uniqueness is the same injectivity: any semifunctor with
$\Ufun\circ\Phiopen=\Phiphase$ must send $f$ to a machine whose
$\Ufun$-image is $\Phiphase(f)$, and $M_f$ is the only one.
\end{proof}

\subsection{Machine embeddings and colimits}
\label{sec.machine_colimits}

Machines at different ports compare along embeddings of ports, and
towers of such embeddings have colimits, computed by completed unions.
We record the infrastructure here; \cref{sec.divisible} uses it.

\begin{definition}[Port embedding]\label{def.port_embedding}
A \emph{port embedding} $u\colon(A,\lambda)\to(A',\lambda')$ is an
isometry $u\colon A\to A'$ satisfying
\[
(u^\dagger)^*\circ\lambda=\lambda'\circ u;
\]
equivalently, by \cref{eq.riesz_identities}, $u$ intertwines the
impedance operators $\riesz_A\circ\lambda$ and
$\riesz_{A'}\circ\lambda'$.
\end{definition}

Port embeddings compose, since transposes and daggers of composites
are composites; and the form pulls back along a port embedding:
$u^*\circ\lambda'\circ u=(u^\dagger u)^*\circ\lambda=\lambda$.

\begin{definition}[Morphisms over port embeddings]\label{def.machine_over}
Let $M=(S,f^+,f^-)\colon(A,\lambda)\to(B,\mu)$ and
$M'=(S',f'^+,f'^-)\colon(A',\lambda')\to(B',\mu')$ be ported machines
and let $u\colon(A,\lambda)\to(A',\lambda')$ and
$v\colon(B,\mu)\to(B',\mu')$ be port embeddings.  A \emph{morphism
$\alpha\colon M\to M'$ over $(u,v)$} is a bounded linear map
$\alpha\colon S\to S'$ satisfying
\[
f'^+\circ\alpha=\bigl(v\oplus(u^\dagger)^*\bigr)\circ f^+
\qquad\text{and}\qquad
\alpha\circ f^-
=f'^-\circ\bigl(\alpha\oplus u\oplus(v^\dagger)^*\bigr).
\]
\end{definition}

Values travel along the embeddings, and covectors along the
transposes $(u^\dagger)^*$ of their coisometries---the Riesz
conjugates of the embeddings themselves \cref{eq.riesz_identities},
in particular isometries.  At $(u,v)=(\id,\id)$ these are the
morphisms of \cref{def.ported_machine}, and morphisms compose over
composite embeddings, since
$\bigl((u'\circ u)^\dagger\bigr)^*=(u'^\dagger)^*\circ(u^\dagger)^*$.

\begin{lemma}[Horizontal composition]\label{lem.horizontal_composition}
Let $\alpha\colon M\to M'$ be a morphism over $(u,v)$ and
$\beta\colon N\to N'$ a morphism over $(v,x)$, with $N\circ M$ and
$N'\circ M'$ defined.  Then $\alpha\oplus\beta$ is a morphism
$N\circ M\to N'\circ M'$ over $(u,x)$.
\end{lemma}

\begin{proof}
Unwind \cref{eq.machine_composition}.  Every exchange is a defining
condition of $\alpha$ or $\beta$ except the middle insertion, where
the port-embedding condition for $v$ gives
$(v^\dagger)^*\bigl(\mu(b)+\zeta\bigr)=\mu'(vb)+(v^\dagger)^*\zeta$
for the value $b$ of $M$ and the offset $\zeta$ of $N$.
\end{proof}

\begin{proposition}[Machine colimits]\label{prop.machine_colimits}
Let $(M_m)_{m\geq n}$ be ported machines
$M_m\colon(A_m,\lambda_m)\to(B_m,\mu_m)$ with morphisms
$\alpha_m\colon M_m\to M_{m+1}$ over port embeddings $(u_m,v_m)$,
such that every $\alpha_m$ is an isometry and the norms
$\lVert f_m^{\pm}\rVert$, $\lVert\lambda_m\rVert$, and
$\lVert\mu_m\rVert$ are bounded uniformly in $m$.  Then the completed
unions of the port and state towers carry a unique ported machine
\[
M_\infty\colon(A_\infty,\lambda_\infty)\to(B_\infty,\mu_\infty)
\]
whose impedance forms and structure maps restrict to those of the
tower; the canonical maps $M_m\to M_\infty$ are morphisms over the
canonical port embeddings; and for every ported machine $N$ with a
compatible family of morphisms $M_m\to N$ over port embeddings whose
state maps are uniformly bounded, exactly one morphism
$M_\infty\to N$ restricts to the family.  We write
$M_\infty=\operatorname{colim}_mM_m$.
\end{proposition}

\begin{proof}
Write $\iota_m\colon A_m\to A_\infty$, $\jmath_m\colon B_m\to
B_\infty$, and $\kappa_m\colon S_m\to S_\infty$ for the embeddings
into the completed unions.

For the impedance forms: the impedance operators
$L_m\coloneqq\riesz_{A_m}\circ\lambda_m$ satisfy
$L_{m+1}u_m=u_mL_m$ (\cref{def.port_embedding}), so the maps
$\iota_mL_m$ are a uniformly bounded compatible family, and
\cref{lem.isometric_colimits} gives a unique bounded
$L_\infty\colon A_\infty\to A_\infty$ with
$L_\infty\iota_m=\iota_mL_m$; it is self-adjoint, being symmetric on
the dense union, and $\lambda_\infty\coloneqq\riesz_{A_\infty}^{-1}
\circ L_\infty$ is a symmetric form for which every $\iota_m$ is a
port embedding.  Likewise $\mu_\infty$.

For the structure maps: the family
$\bigl(\jmath_m\oplus(\iota_m^\dagger)^*\bigr)\circ f_m^+$ is
compatible by the defining conditions of the $\alpha_m$ and uniformly
bounded, so it extends uniquely to
$f^+_\infty\colon S_\infty\to B_\infty\oplus A_\infty^*$
(\cref{lem.isometric_colimits}).  For the update: direct sums commute
with completed unions, and the covector legs $(v_m^\dagger)^*$ are
the Riesz conjugates of the $v_m$, so the completed union of
$S_m\oplus A_m\oplus B_m^*$ along
$\alpha_m\oplus u_m\oplus(v_m^\dagger)^*$ is
$S_\infty\oplus A_\infty\oplus B_\infty^*$, and the family
$\kappa_m\circ f_m^-$ extends uniquely to $f^-_\infty$.  The two
conditions of \cref{def.machine_over} for the canonical maps hold
because both sides are bounded and agree on the dense unions; the
same argument gives the uniqueness of $M_\infty$.

For the universal property: \cref{lem.isometric_colimits} supplies
unique bounded port and state maps restricting to the given family;
the port maps are isometries, being isometric on dense subspaces, and
are port embeddings, and the state map satisfies the conditions of
\cref{def.machine_over}, in each case because both sides of the
required identity are bounded and agree on a dense union.
\end{proof}

\section{The cell and its chains}
\label{sec.cell}

Fix an impedance $Z\in\R_{>0}$.

\begin{definition}[The $Z$-port and the cell]\label{def.cell}
The \emph{$Z$-port} is
\[
\defineTerm{Zport}\coloneqq\bigl(\R,\ Z\,\riesz_\R^{-1}\bigr),
\]
the line with impedance form $a\mapsto Z\langle a,\blank\rangle$.  The
\emph{cell} is the morphism
\[
\defineTerm{cellm}_Z\colon \Zport\longrightarrow \Zport
\quad\text{in }\Open
\]
with parameter $(\R,\ \riesz_\R/Z)$, output map $q\mapsto q$, and
potential
\[
U(q,a)\coloneqq\frac Z2(a-q)^2.
\]
\end{definition}

The cell is a unit mass bound by a unit bond to its input port: for a
cell duration $h\in\R_{>0}$, a mass of size $Zh$ has reaction
$\riesz/(Zh)$ and a bond of stiffness $Z/h$ has potential
$\frac{Z}{2h}(a-q)^2$, and executing one $\pc$-tick per duration $h$
multiplies both rates by $h$ (\cref{sec.hilbert_semantics}), leaving
the tuned data of \cref{def.cell}.  The potential is a continuous
quadratic form on $\R\oplus\R$, the reaction is symmetric, the output
map is bounded linear, and
$d_AU|_{(0,a)}=Z\langle a,\blank\rangle$, so the impedance condition
\cref{eq.impedance_condition} holds at $\Zport$: the cell is a morphism
of $\Open$.  The output port carries the bond of the \emph{next} cell,
which is why the cell has a bond at its input only.

\begin{definition}[Chains]\label{def.chain}
For $K\in\Z_{\geq1}$, the \emph{$K$-chain} is the $K$th power
\[
\defineTerm{Chain}_K\coloneqq\cellm_Z^{\circ K}
\colon \Zport\longrightarrow \Zport
\quad\text{in }\Open.
\]
\end{definition}

\begin{proposition}[Chain data]\label{prop.chain_data}
The $K$-chain is the open system with parameter
$\bigl(\ell^2(\{1,\ldots,K\}),\ \riesz/Z\bigr)$, output map
$q\mapsto q_K$, and potential
\begin{equation}\label{eq.chain_potential}
U_K(q,a)
=
\frac Z2
\Bigl(
(a-q_1)^2+\sum_{i=1}^{K-1}(q_i-q_{i+1})^2
\Bigr).
\end{equation}
\end{proposition}

\begin{proof}
Induction on $K$.  The case $K=1$ is \cref{def.cell}.  For the
inductive step, $\Chain_{K+1}=\cellm_Z\circ\Chain_K$, and
\cref{eq.open_composition} gives parameter
$(\ell^2(\{1,\ldots,K\})\oplus\R,\ \riesz/Z)$
(\cref{prop.rhilb_monoidal}), output map the last coordinate, and
potential
\[
U_K(q,a)+U\bigl(q',\ q_K\bigr)
=
U_K(q,a)+\frac Z2(q_K-q')^2,
\]
which is \cref{eq.chain_potential} for $K+1$ after writing
$q_{K+1}\coloneqq q'$.
\end{proof}

\begin{corollary}[The chain semigroup]\label{cor.chain_semigroup}
$\Chain_{K+K'}=\Chain_{K'}\circ\Chain_K$ for all $K,K'\in\Z_{\geq1}$:
the assignment $K\mapsto\Chain_K$ is a homomorphism of semigroups
$(\Z_{\geq1},+)\to\Open(\Zport,\Zport)$.
\end{corollary}

\begin{proof}
Powers in a semicategory satisfy
$x^{\circ(K+K')}=x^{\circ K'}\circ x^{\circ K}$ by associativity.
\end{proof}

\begin{definition}[Chain machine]\label{def.chain_machine}
The \emph{$K$-chain machine} is
\[
\defineTerm{Tm}_Z(K)\coloneqq\Phiopen(\Chain_K)
\colon \Zport\longrightarrow \Zport
\quad\text{in }\Lin.
\]
\end{definition}

\begin{corollary}[Chain machines compose]\label{cor.chain_machines_compose}
$\Tm_Z(K{+}K')=\Tm_Z(K')\circ\Tm_Z(K)$ for all $K,K'\in\Z_{\geq1}$:
the assignment $K\mapsto\Tm_Z(K)$ is a homomorphism of semigroups.
\end{corollary}

\begin{proof}
Apply the semifunctor $\Phiopen$ (\cref{thm.open_semantics}) to
\cref{cor.chain_semigroup}.
\end{proof}

\section{The characteristic flow}
\label{sec.characteristic}

The chain machine transports wave packets exactly.  This section
defines the packets of a driven chain trajectory and proves the
transport law; the packets are staggered differences of the extended
displacement field, so they live on trajectories rather than on
states, and the statements of this section are trajectory statements.

\begin{definition}[Driven trajectory]\label{def.trajectory}
Let $M=(S,f^+,f^-)\colon(A,\lambda)\to(B,\mu)$ be a ported machine.  A
\emph{driven trajectory} of $M$ is a triple of sequences
\[
s\colon\N\to S,
\qquad
a\colon\N\to A,
\qquad
\eta\colon\N\to B^*
\]
with $s(t+1)=f^-\bigl(s(t),a(t),\eta(t)\bigr)$ for every $t\in\N$; the
\emph{emissions} of the trajectory are the readouts $f^+(s(t))$.
\end{definition}

Driven trajectories are the $(\N,\mathrm{succ})$-shaped elements of a
closed system in $\pc$: closing the machine with an environment that
plays the streams leaves a discrete dynamical system, and a trajectory
is a morphism into it from the successor system.

\begin{definition}[Stream environment]\label{def.stream_environment}
Let $(A,\lambda)$ and $(B,\mu)$ be ports and let $a\colon\N\to A$ and
$\eta\colon\N\to B^*$ be sequences, the \emph{streams}.  The
\emph{stream environment}
\[
E_{a,\eta}\in
\pc\bigl(\cotof B\otimes[\cotof A,\yon],\ \yon\bigr)
\]
is the $1$-cell with carrier $\N$ whose direction map at $t\in\N$
returns, at every position of the domain, the direction
$\bigl(\eta(t),a(t)\bigr)$ and the next state $t+1$.
\end{definition}

\begin{proposition}[Trajectories are elements]
\label{prop.trajectories_are_elements}
Let $M$ be a ported machine, let $(a,\eta)$ be streams for its ports,
and write $(\N,\mathrm{succ})\in\pc(\yon,\yon)$ for the successor
dynamical system.  The assignment
$s\mapsto\bigl(t\mapsto(s(t),t)\bigr)$ is a bijection between driven
trajectories of $M$ with streams $(a,\eta)$ and $2$-cells
\[
(\N,\mathrm{succ})
\Longrightarrow
E_{a,\eta}\circ\Ufun(M)
\quad\text{in }\pc(\yon,\yon)
\]
whose clock component at $0$ is $0$.
\end{proposition}

\begin{proof}
Unwinding the Dirichlet composite of coalgebras
(\cref{sec.background_categories}) as in \cref{thm.open_semantics}:
the composite $E_{a,\eta}\circ\Ufun(M)$ is the dynamical system with
state set $\lvert S\rvert\times\N$ in which one tick presents the
position of $\Ufun(M)$ at $s$, receives the direction
$\bigl(\eta(t),a(t)\bigr)$ from the environment, and updates to
$\bigl(f^-(s,a(t),\eta(t)),\ t+1\bigr)$
(\cref{def.ufun,def.stream_environment}).  A $2$-cell from
$(\N,\mathrm{succ})$ is a function $t\mapsto\bigl(s(t),c(t)\bigr)$
commuting with the two transitions
(\cref{sec.background_categories}).  The clock component gives
$c(t+1)=c(t)+1$, so $c=\id_\N$ exactly when $c(0)=0$, and the state
component is then the condition
$s(t+1)=f^-\bigl(s(t),a(t),\eta(t)\bigr)$ of \cref{def.trajectory}.
\end{proof}

\begin{definition}[Chain trajectory notation]\label{def.chain_notation}
Let $K\in\Z_{\geq1}$ and let $\bigl((q,\xi),a,\eta\bigr)$ be a driven
trajectory of the chain machine $\Tm_Z(K)$
(\cref{def.trajectory,def.chain_machine}): at each tick $t\in\N$, the
state is $\bigl(q(t),\xi(t)\bigr)\in T^*\ell^2(\{1,\ldots,K\})$, the
input is $a(t)\in\R$, and the incoming covector is $\eta(t)\in\R^*$.
We write
\[
\tilde q(t)\coloneqq q(t)+\riesz\bigl(\xi(t)\bigr)/Z,
\qquad
p(t)\coloneqq\riesz\bigl(\xi(t)\bigr)/Z,
\qquad
e(t)\coloneqq\riesz\bigl(\eta(t)\bigr)
\]
for the presented position, the velocity, and the force value, and
\[
\dot q(t)\coloneqq q(t+1)-q(t),
\qquad
\ddot q(t)\coloneqq q(t+1)-2q(t)+q(t-1)
\quad(t\geq1),
\]
for the discrete velocity and the discrete acceleration.
\end{definition}

Unwinding $\Tm_Z(K)=\Phiopen(\Chain_K)$ by
\cref{def.machine_of_system,prop.chain_data}: the emissions at
$\bigl(q(t),\xi(t)\bigr)$ are the value $\tilde q_K(t)$ and the offset
$-Z\,\riesz_\R^{-1}\bigl(\tilde q_1(t)\bigr)$, and the trajectory
condition of \cref{def.trajectory} reads
\begin{equation}\label{eq.chain_update}
q(t+1)=\tilde q(t),
\qquad
\xi(t+1)=\xi(t)-(\outp{\Chain_K})^*\eta(t)
-d_QU_K|_{(\tilde q(t),a(t))}.
\end{equation}
In particular the position component gives
$q(t+1)=q(t)+p(t)$, i.e.
\begin{equation}\label{eq.velocity}
\dot q(t)=p(t).
\end{equation}

\begin{definition}[Extended displacement]\label{def.extended_displacement}
Let $\bigl((q,\xi),a,\eta\bigr)$ be a driven trajectory of
$\Tm_Z(K)$, with notation as in \cref{def.chain_notation}.  Its
\emph{extended displacement} is the family
$\bar q_i(t)$, defined for $1\leq i\leq K$ at all $t\in\N$ and for
$i\in\{0,K{+}1\}$ at $t\geq1$, by
\[
\bar q_i(t)\coloneqq q_i(t),
\qquad
\bar q_0(t)\coloneqq a(t-1),
\qquad
\bar q_{K+1}(t)\coloneqq \tilde q_K(t-1)-e(t-1)/Z:
\]
the input plays the presented position of a virtual site $0$, and the
incoming covector deposits a virtual site $K+1$.
\end{definition}

\begin{lemma}[Extended recurrence]\label{lem.extended_recurrence}
Let $\bigl((q,\xi),a,\eta\bigr)$ be a driven trajectory of
$\Tm_Z(K)$, with extended displacement $\bar q$
(\cref{def.extended_displacement}) and discrete acceleration $\ddot q$
(\cref{def.chain_notation}).  For $1\leq i\leq K$ and $t\geq1$,
\begin{equation}\label{eq.extended_recurrence}
\ddot q_i(t)=\bar q_{i-1}(t)-2\bar q_i(t)+\bar q_{i+1}(t).
\end{equation}
\end{lemma}

\begin{proof}
By \cref{eq.velocity}, $\ddot q_i(t)=p_i(t)-p_i(t-1)$, and the
momentum component of \cref{eq.chain_update} at tick $t-1$ computes
it.  Differentiating the chain potential \cref{eq.chain_potential},
the components of $\riesz\bigl(d_QU_K|_{(\tilde q,a)}\bigr)/Z$ are
$(\tilde q_1-a)+(\tilde q_1-\tilde q_2)$ at $i=1$,
$\ 2\tilde q_i-\tilde q_{i-1}-\tilde q_{i+1}$ at $1<i<K$, and
$\tilde q_K-\tilde q_{K-1}$ at $i=K$ (at $K=1$, the single component
is $\tilde q_1-a$); and $(\outp{\Chain_K})^*\eta=\eta$ in the $K$th
component.  Using $\tilde q_j(t-1)=q_j(t)$ throughout
\cref{eq.velocity}: for $1<i<K$,
\[
p_i(t)-p_i(t-1)
=-\bigl(2\tilde q_i-\tilde q_{i-1}-\tilde q_{i+1}\bigr)(t-1)
=q_{i-1}(t)-2q_i(t)+q_{i+1}(t).
\]
For $i=1$ (and $K\geq2$):
\[
p_1(t)-p_1(t-1)
=\bigl(a(t-1)-\tilde q_1(t-1)\bigr)
+\bigl(\tilde q_2(t-1)-\tilde q_1(t-1)\bigr)
=\bigl(\bar q_0(t)-q_1(t)\bigr)+\bigl(q_2(t)-q_1(t)\bigr).
\]
For $i=K$ (and $K\geq2$):
\[
p_K(t)-p_K(t-1)
=\bigl(\tilde q_{K-1}(t-1)-\tilde q_K(t-1)\bigr)-e(t-1)/Z
=\bigl(q_{K-1}(t)-q_K(t)\bigr)
+\bigl(\bar q_{K+1}(t)-q_K(t)\bigr),
\]
since $\bar q_{K+1}(t)=q_K(t)-e(t-1)/Z$.  For $K=1$ the two boundary
corrections combine:
$p_1(t)-p_1(t-1)
=\bigl(\bar q_0(t)-q_1(t)\bigr)+\bigl(\bar q_2(t)-q_1(t)\bigr)$.
\end{proof}

\begin{definition}[Packets]\label{def.packets}
Let $\bigl((q,\xi),a,\eta\bigr)$ be a driven trajectory of
$\Tm_Z(K)$, with extended displacement $\bar q$
(\cref{def.extended_displacement}).  The \emph{packets} of the
trajectory are
\[
w^{\rightarrow}_i(t)
\coloneqq
\tfrac12\bigl(\bar q_i(t+1)-\bar q_{i+1}(t)\bigr)
\quad(0\leq i\leq K),
\qquad
w^{\leftarrow}_i(t)
\coloneqq
\tfrac12\bigl(\bar q_i(t+1)-\bar q_{i-1}(t)\bigr)
\quad(1\leq i\leq K{+}1),
\]
defined whenever the entries are.  The packets in flight
$\bigl(w^{\rightarrow}_i(t),w^{\leftarrow}_i(t)\bigr)_{1\leq i\leq K}$
form the \emph{wave chart} of the trajectory at time $t$.  The
\emph{in-streams} are the boundary packets
\begin{equation}\label{eq.in_streams}
w^{\rightarrow}_0(t)
=\tfrac12\bigl(a(t)-q_1(t)\bigr),
\qquad
w^{\leftarrow}_{K+1}(t)
=\tfrac12\bigl(p_K(t)-e(t)/Z\bigr),
\end{equation}
defined for all $t\in\N$; the \emph{out-streams} are
\begin{equation}\label{eq.out_streams}
w^{\rightarrow}_K(t)
=\tfrac12\bigl(\tilde q_K(t)-\tilde q_K(t-1)+e(t-1)/Z\bigr),
\qquad
w^{\leftarrow}_1(t)
=\tfrac12\bigl(\tilde q_1(t)-a(t-1)\bigr),
\end{equation}
defined for $t\geq1$.
\end{definition}

The identities \cref{eq.in_streams,eq.out_streams} are immediate from
\cref{def.extended_displacement}, using \cref{eq.velocity} for the
second in-stream:
$\bar q_{K+1}(t+1)-q_K(t)=\tilde q_K(t)-e(t)/Z-q_K(t)=p_K(t)-e(t)/Z$.
Each is a bounded linear expression in the port traffic and the state:
the in-streams in the current input, covector, and state; the
out-streams in the current and previous emissions and covector.

\begin{theorem}[Characteristic flow]\label{thm.characteristic_flow}
Let $\bigl((q,\xi),a,\eta\bigr)$ be a driven trajectory of
$\Tm_Z(K)$, with packets $w^{\rightarrow},w^{\leftarrow}$
(\cref{def.packets}).  The trajectory satisfies the transport laws
\begin{equation}\label{eq.transport}
w^{\rightarrow}_i(t+1)=w^{\rightarrow}_{i-1}(t),
\qquad
w^{\leftarrow}_i(t+1)=w^{\leftarrow}_{i+1}(t)
\qquad(1\leq i\leq K,\ t\geq0).
\end{equation}
Consequently, for $t\geq K$, the out-streams are the in-streams
delayed by exactly $K$ ticks,
\begin{equation}\label{eq.exact_delay}
w^{\rightarrow}_K(t)=w^{\rightarrow}_0(t-K),
\qquad
w^{\leftarrow}_1(t)=w^{\leftarrow}_{K+1}(t-K),
\end{equation}
and the packets in flight record the in-streams:
\begin{equation}\label{eq.chart_records}
w^{\rightarrow}_i(t)=w^{\rightarrow}_0(t-i),
\qquad
w^{\leftarrow}_i(t)=w^{\leftarrow}_{K+1}\bigl(t-(K{+}1{-}i)\bigr)
\qquad(1\leq i\leq K).
\end{equation}
\end{theorem}

\begin{proof}
For the first law of \cref{eq.transport}, let $1\leq i\leq K$ and
$t\geq0$; every extended displacement appearing below is defined.  By
\cref{def.packets},
\[
2\,w^{\rightarrow}_i(t+1)-2\,w^{\rightarrow}_{i-1}(t)
=\bar q_i(t+2)-\bar q_{i+1}(t+1)-\bar q_{i-1}(t+1)+\bar q_i(t).
\]
Since $1\leq i\leq K$, $\bar q_i=q_i$, so
$\bar q_i(t+2)+\bar q_i(t)=\ddot q_i(t+1)+2q_i(t+1)$, and the
right-hand side becomes
\[
\ddot q_i(t+1)+2\bar q_i(t+1)-\bar q_{i+1}(t+1)-\bar q_{i-1}(t+1),
\]
which vanishes by the extended recurrence
\cref{eq.extended_recurrence} at $(i,t+1)$.  The second law is the
same computation with $i+1$ and $i-1$ exchanged.  Iterating
\cref{eq.transport} gives \cref{eq.exact_delay,eq.chart_records}: each
application lowers the time by one and moves the site by one toward
the entry port, and the iterates stay in range until they reach the
in-streams \cref{eq.in_streams}, which are defined at all times.
\end{proof}

\section{Queues and strings}
\label{sec.strings}

The state that \cref{thm.characteristic_flow} exhibits---packets in
flight, entering at one port and exiting at the other after a fixed
number of ticks---is a machine in its own right.  This section defines
it, proves it terminal among realizations of exact delay, and shows
that chains are strings in wave variables.

\subsection{Queues}
\label{sec.queues_sub}

Fix a real Hilbert space $B\in\Hilb$ and an integer $N\geq1$.  Define
\[
\defineTerm{Qu}_N(B)\coloneqq B^{\oplus N}
\]
and bounded linear maps, for every $x=(b_0,\ldots,b_{N-1})\in \Qu_N(B)$
and $b\in B$,
\begin{equation}\label{eq.queue_structure}
\begin{aligned}
S_Nx&\coloneqq(b_1,\ldots,b_{N-1},0),\\
J_Nb&\coloneqq(0,\ldots,0,b),
\end{aligned}
\qquad
\begin{aligned}
P_Nx&\coloneqq b_0,\\
U_N(x,b)&\coloneqq S_Nx+J_Nb.
\end{aligned}
\end{equation}
The \emph{queue step} is
\begin{equation}\label{eq.finite_queue_step}
\step_N\colon \Qu_N(B)\oplus B\longrightarrow B\oplus \Qu_N(B),
\qquad
\step_N(x,b)\coloneqq\bigl(P_Nx,\,U_N(x,b)\bigr).
\end{equation}

\begin{definition}[Exact-delay realization]
\label{def.delay_realization}
Fix $B\in\Hilb$ and an integer $N\geq1$.  A \emph{bounded linear
realization of exact $N$-step delay on $B$} consists of a Hilbert space
$X\in\Hilb$ and bounded maps
\[
T_X\colon X\to X,
\qquad
G_X\colon B\to X,
\qquad
R_X\colon X\to B
\]
satisfying
\begin{equation}\label{eq.delay_laws}
R_XT_X^N=0,
\qquad
R_XT_X^kG_X
=
\begin{cases}
0,&0\leq k<N-1,\\
\id_B,&k=N-1.
\end{cases}
\end{equation}
A morphism $f\colon(X,T_X,G_X,R_X)\to(Y,T_Y,G_Y,R_Y)$ is a bounded
linear map $f\colon X\to Y$ satisfying
\[
fT_X=T_Yf,
\qquad
fG_X=G_Y,
\qquad
R_X=R_Yf.
\]
Write $\defineTerm{DelayCat}_N(B)$ for the resulting category.
\end{definition}

\begin{theorem}[Queue representation]\label{thm.queue_representation}
Fix $B\in\Hilb$ and an integer $N\geq1$.

\begin{enumerate}
\item\label{item.step_unitary} The queue step $\step_N$ is unitary.

\item\label{item.terminal} The realization
$\defineTerm{Qterm}_N(B)\coloneqq\bigl(\Qu_N(B),S_N,J_N,P_N\bigr)$
is terminal in $\DelayCat_N(B)$.

\item\label{item.terminal_map} For every object
$M=(X,T_X,G_X,R_X)\in\DelayCat_N(B)$, the terminal map is
\[
\Theta_M\colon X\longrightarrow \Qu_N(B),
\qquad
\Theta_M(x)\coloneqq
\bigl(R_Xx,R_XT_Xx,\ldots,R_XT_X^{N-1}x\bigr).
\]
Its underlying map in $\Hilb$ is a split epimorphism.  If
\begin{equation}\label{eq.lossless_observability}
\lVert x\rVert_X^2
=
\sum_{k=0}^{N-1}\lVert R_XT_X^kx\rVert_B^2
\qquad(x\in X),
\end{equation}
then $\Theta_M$ is a unitary isomorphism in $\DelayCat_N(B)$.
\end{enumerate}
\end{theorem}

\begin{proof}
The map $\step_N$ is the coordinate permutation
\[
\bigl((b_0,\ldots,b_{N-1}),b\bigr)
\longmapsto
\bigl(b_0,(b_1,\ldots,b_{N-1},b)\bigr),
\]
with the first component designated as output and the last $N$ components
designated as the new state; hence it is unitary.

The maps in \cref{eq.queue_structure} satisfy the delay laws
\cref{eq.delay_laws}.  Fix an object
$M=(X,T_X,G_X,R_X)\in\DelayCat_N(B)$.  Its delay laws give
\[
P_N\Theta_M=R_X,
\qquad
\Theta_MT_X=S_N\Theta_M,
\qquad
\Theta_MG_X=J_N.
\]
Thus $\Theta_M$ is a morphism to $\Qterm_N(B)$.  Its $k$th component
must be $R_XT_X^k$, so it is the unique such morphism.

A section of $\Theta_M$ is
\[
L_M\colon \Qu_N(B)\longrightarrow X,
\qquad
L_M(b_0,\ldots,b_{N-1})
\coloneqq
\sum_{k=0}^{N-1}T_X^{N-1-k}G_Xb_k.
\]
Indeed, \cref{eq.delay_laws} gives
$\Theta_ML_M=\id_{\Qu_N(B)}$.  Finally,
\cref{eq.lossless_observability} says exactly that $\Theta_M$ is isometric;
the split epimorphism is then unitary.
\end{proof}

\subsection{Strings}
\label{sec.strings_sub}

\begin{definition}[Strings]\label{def.string}
For $B\in\Hilb$ and $K\in\Z_{\geq1}$, the \emph{$K$-string on $B$} is
the ported machine
\[
\defineTerm{Sm}_B(K)\colon(B,0)\longrightarrow(B,0)
\quad\text{in }\Lin
\]
with state space $\Qu_K(B)\oplus \Qu_K(B^*)$ and
\[
\Sm_B(K)^+(x,y)\coloneqq\bigl(P_Kx,\ P_Ky\bigr),
\qquad
\Sm_B(K)^-\bigl((x,y),a,\theta\bigr)
\coloneqq\bigl(U_K(x,a),\ U_K(y,\theta)\bigr),
\]
where $P_K$ and $U_K$ are the queue readout and update
\cref{eq.queue_structure}.
\end{definition}

The rightward queue $x$ receives the input value and its head is the
emitted value; the leftward queue $y$ receives the incoming covector
and its head is the emitted offset.  The port impedance form is zero:
the string returns nothing instantaneous along its input, only the
head of its leftward queue.

\begin{proposition}[Strings compose]\label{prop.string_composition}
For $K,K'\in\Z_{\geq1}$, the concatenation
\[
\chi_{K,K'}\bigl((x,y),(x',y')\bigr)
\coloneqq
\bigl((x',x),\ (y,y')\bigr)
\]
is a unitary $2$-cell
$\chi_{K,K'}\colon
\Sm_B(K')\circ\Sm_B(K)\To{\cong}\Sm_B(K{+}K')$
in $\Lin\bigl((B,0),(B,0)\bigr)$, and the family $(\chi_{K,K'})$
satisfies the associativity coherence: the assignment
$K\mapsto\Sm_B(K)$ is a strong semigroupal representation of
$(\Z_{\geq1},+)$ in
$\bigl(\Lin((B,0),(B,0)),\circ\bigr)$.
\end{proposition}

\begin{proof}
The map $\chi_{K,K'}$ permutes direct summands, hence is unitary.  We
check the machine commutations.  Readouts: the composite's readout
\cref{eq.machine_composition} is
$\bigl(P_{K'}x',\ P_Ky\bigr)$, and the concatenations $(x',x)$ and
$(y,y')$ have heads $x'_0$ and $y_0$.  Updates: in the composite, the
first string receives the input $a$ and, the middle impedance form
being zero, the incoming covector $P_{K'}y'$; the second receives the
value $P_Kx$ and the covector $\theta$.  So the composite update is
\[
\Bigl(
\bigl(U_K(x,a),\ U_K(y,P_{K'}y')\bigr),\
\bigl(U_{K'}(x',P_Kx),\ U_{K'}(y',\theta)\bigr)
\Bigr).
\]
Applying $\chi_{K,K'}$ gives
$\bigl((U_{K'}(x',P_Kx),U_K(x,a)),\ (U_K(y,P_{K'}y'),U_{K'}(y',\theta))\bigr)$,
and by \cref{eq.queue_structure} this is exactly
$\bigl(U_{K+K'}((x',x),a),\ U_{K+K'}((y,y'),\theta)\bigr)$: shifting
the concatenation moves the head of $x$ into the tail of $x'$ and
inserts $a$ at the far end, and symmetrically for the leftward queue.
Coherence: both composites
$\chi\circ(\chi\circ\id)$ and $\chi\circ(\id\circ\chi)$ on a triple
are the same permutation of summands, by associativity of
concatenation.
\end{proof}

\begin{proposition}[Strings are functorial]\label{prop.string_functorial}
Let $K\in\Z_{\geq1}$.  An isometry $w\colon B\to B'$ induces the
morphism
\[
\Sm_w(K)\coloneqq \Qu_K(w)\oplus \Qu_K\bigl((w^\dagger)^*\bigr)
\colon\Sm_B(K)\longrightarrow\Sm_{B'}(K)
\]
over the port embedding $(w,w)$ of zero-impedance ports
(\cref{def.machine_over}).  This makes $\Sm_{(-)}(K)$ a functor from
Hilbert spaces and isometries to ported machines and morphisms over
port embeddings; it preserves tower colimits
(\cref{prop.machine_colimits}); and the concatenation $2$-cells
$\chi_{K,K'}$ of \cref{prop.string_composition} are natural in $B$.
\end{proposition}

\begin{proof}
The queue maps \cref{eq.queue_structure} permute, insert, and read
coordinates, and $\Qu_K(w)$ acts slotwise, so the two conditions of
\cref{def.machine_over} hold componentwise on the value channel; on
the covector channel they are the Riesz conjugates of the same
identities, $(w^\dagger)^*$ being the Riesz conjugate of $w$
\cref{eq.riesz_identities}, in particular an isometry.
Functoriality is $\Qu_K$ applied to composites.

For tower colimits: given a tower of isometries $(w_m)$ with
completed union $B_\infty$, the machines $\Sm_{B_m}(K)$ and morphisms
$\Sm_{w_m}(K)$ satisfy the hypotheses of
\cref{prop.machine_colimits}---the state maps are isometries and the
structure maps are contractions---and finite direct sums commute with
completed unions, so the state tower completes to
$\Qu_K(B_\infty)\oplus \Qu_K(B_\infty^*)$, the covector channel
through the Riesz conjugates.  The structure maps of
$\Sm_{B_\infty}(K)$ restrict on the dense unions to those of the
tower, hence coincide with those of the colimit machine.

Naturality: $\chi_{K,K'}$ and $\chi'_{K,K'}$ permute direct summands
and $\Sm_w$ acts slotwise, so the two composites around the
naturality square are equal.
\end{proof}

\begin{proposition}[String trajectories]\label{prop.string_trajectory}
Let $\bigl((x,y),a,\theta\bigr)$ be a driven trajectory of
$\Sm_B(K)$.  For $0\leq j<K$ and $t\geq K-j$,
\[
x(t)_j=a\bigl(t-(K-j)\bigr),
\qquad
y(t)_j=\theta\bigl(t-(K-j)\bigr);
\]
in particular the emitted value at time $t\geq K$ is $a(t-K)$ and the
emitted offset is $\theta(t-K)$.
\end{proposition}

\begin{proof}
The update inserts $a(t)$ at slot $K-1$ and shifts every slot down by
one per tick, so slot $j$ at time $t$ holds the input from time
$t-(K-j)$ once that input exists; likewise for the leftward queue.
The emissions are the heads $x(t)_0$ and $y(t)_0$.
\end{proof}

\subsection{Chains are strings in wave variables}
\label{sec.chain_string}

\begin{corollary}[Chain--string comparison]\label{cor.chain_string}
Let $\bigl((q,\xi),a,\eta\bigr)$ be a driven trajectory of
$\Tm_Z(K)$, with packets $w^{\rightarrow},w^{\leftarrow}$
(\cref{def.packets}), and let $\bigl((x,y),\beta,\vartheta\bigr)$ be
the driven trajectory of the string $\Sm_\R(K)$ with any initial
state and the in-streams of the chain as driving:
\[
\beta(t)\coloneqq w^{\rightarrow}_0(t),
\qquad
\vartheta(t)\coloneqq\riesz_\R^{-1}\bigl(w^{\leftarrow}_{K+1}(t)\bigr).
\]
Then for every $t\geq K$ and $0\leq j<K$,
\[
x(t)_j=w^{\rightarrow}_{K-j}(t),
\qquad
\riesz_\R\bigl(y(t)_j\bigr)=w^{\leftarrow}_{j+1}(t):
\]
from the first transit onward, the state of the string is the wave
chart of the chain, and the out-streams of the two machines coincide.
\end{corollary}

\begin{proof}
By \cref{prop.string_trajectory},
$x(t)_j=\beta(t-(K-j))=w^{\rightarrow}_0(t-(K-j))$, which equals
$w^{\rightarrow}_{K-j}(t)$ by \cref{eq.chart_records}; and
$\riesz_\R(y(t)_j)=w^{\leftarrow}_{K+1}(t-(K-j))
=w^{\leftarrow}_{(j+1)}(t)$, again by \cref{eq.chart_records} with
$K{+}1{-}(j{+}1)=K-j$.  At $j=0$ these are the out-streams
\cref{eq.out_streams} of the chain and the emissions of the string.
\end{proof}

\section{Subdivision and completion}
\label{sec.queues}

The strings of \cref{sec.strings} exist at every resolution, and
resolutions compare: one coarse packet is two fine packets, and one
coarse tick is two fine ticks.  This section weights the packets,
proves the subdivision theorem, and completes the dyadic tower.

\subsection{Packets and subdivision}
\label{sec.subdivision}

\begin{definition}[Cell packet]\label{def.cell_packet}
For $Z,h\in\R_{>0}$, the \emph{cell packet} is
\[
\defineTerm{CZ}(h)\coloneqq\R[Zh]
\]
(\cref{sec.background_hilbert}): the real line with inner product
$\langle a,b\rangle_{\CZ(h)}=Zh\,ab$.
\end{definition}

For a real Hilbert space $C$, integers $L,r\geq1$, and the queue
step $\step_L$ on $C$ \cref{eq.finite_queue_step}, define the
\emph{$r$-fold block step}
\[
\step_L^{[r]}\colon
\Qu_L(C)\oplus C^{\oplus r}
\longrightarrow
C^{\oplus r}\oplus \Qu_L(C)
\]
as follows.  Given $x_0\in \Qu_L(C)$ and
$(c_0,\ldots,c_{r-1})\in C^{\oplus r}$, define
\[
(y_j,x_{j+1})\coloneqq\step_L(x_j,c_j)
\quad(0\leq j<r),
\qquad
\step_L^{[r]}(x_0,c_0,\ldots,c_{r-1})
\coloneqq(y_0,\ldots,y_{r-1},x_r).
\]

\begin{theorem}[Queue subdivision]\label{thm.queue_subdivision}
Fix real Hilbert spaces $B,C\in\Hilb$, integers $N,L,r\geq1$, and an
isometric embedding
$\sigma\colon B\hookrightarrow C^{\oplus r}$.
Let
$\lambda_{N,r}\colon
\Qu_N(C^{\oplus r})\To{\cong}\Qu_{rN}(C)$
be the unitary that removes parentheses while preserving order, and define
\[
\rho_{N,r,\sigma}
\coloneqq
\lambda_{N,r}\circ \Qu_N(\sigma)
\colon
\Qu_N(B)\lhook\joinrel\longrightarrow \Qu_{rN}(C).
\]
Then, under the identification $\Qu_L(B)=B^{\oplus L}$, the following
square commutes:
\begin{equation}\label{eq.queue_subdivision_square}
\begin{CD}
\Qu_N(B)\oplus \Qu_L(B)
@>{\step_N^{[L]}}>>
\Qu_L(B)\oplus \Qu_N(B)\\
@V{\rho_{N,r,\sigma}\oplus\rho_{L,r,\sigma}}VV
@VV{\rho_{L,r,\sigma}\oplus\rho_{N,r,\sigma}}V\\
\Qu_{rN}(C)\oplus \Qu_{rL}(C)
@>>{\step_{rN}^{[rL]}}>
\Qu_{rL}(C)\oplus \Qu_{rN}(C).
\end{CD}
\end{equation}
In particular, at $L=1$ and zero input, for the queue shifts
$S_N,S_{rN}$ of \cref{eq.queue_structure},
\begin{equation}\label{eq.queue_dilated_equivariance}
\rho_{N,r,\sigma}S_N
=
S_{rN}^{\,r}\rho_{N,r,\sigma},
\end{equation}
so
\begin{equation}\label{eq.queue_dyn_morphism}
\rho_{N,r,\sigma}\colon
\bigl(\Qu_N(B),S_N\bigr)
\longrightarrow
(r\cdot)^*\bigl(\Qu_{rN}(C),S_{rN}\bigr)
\quad\text{in }\Dyn,
\end{equation}
the time change along multiplication by $r$ on $\N$
(\cref{sec.background_hilbert}).
\end{theorem}

\begin{proof}
The map $\rho_{N,r,\sigma}$ replaces each coarse packet by its ordered
$r$-packet subdivision.  The top of \cref{eq.queue_subdivision_square}
removes the first $L$ coarse packets and appends the $L$ coarse
inputs.  The bottom removes their $rL$ refined packets and appends the
$rL$ components of the refined inputs.  Both routes are therefore the
same coordinate reassociation.  Setting $L=1$ and the input equal to
zero gives
\cref{eq.queue_dilated_equivariance,eq.queue_dyn_morphism}.
\end{proof}

\begin{corollary}[Dyadic subdivision]
\label{cor.dyadic_subdivision}
Fix $Z,d\in\R_{>0}$.  For every $n\in\N$, define
\begin{equation}\label{eq.dyadic_data}
K_n\coloneqq2^n,
\qquad
h_n\coloneqq d/K_n,
\qquad
C_n\coloneqq \CZ(h_n),
\qquad
H_n\coloneqq \Qu_{K_n}(C_n).
\end{equation}
The packet map
$\sigma_n\colon C_n\hookrightarrow C_{n+1}^{\oplus2}$,
$\sigma_n(b)\coloneqq(b,b)$,
is an isometric embedding, so the state refinement
\begin{equation}\label{eq.dyadic_state_refinement}
\rho_n
\coloneqq
\lambda_{K_n,2}\circ \Qu_{K_n}(\sigma_n)
\colon H_n\lhook\joinrel\longrightarrow H_{n+1}
\end{equation}
is an isometric embedding satisfying
$\rho_nS_{K_n}=S_{K_{n+1}}^2\rho_n$, and the input-output refinement is
the instance of \cref{eq.queue_subdivision_square} determined by
$B=C_n$, $C=C_{n+1}$, $N=K_n$, $L=1$, $r=2$, and $\sigma=\sigma_n$.
\end{corollary}

\begin{proof}
For every $b\in C_n$,
$\lVert\sigma_n(b)\rVert^2
=2Zh_{n+1}b^2
=Zh_nb^2
=\lVert b\rVert^2$.
The result is now \cref{thm.queue_subdivision}.
\end{proof}

\subsection{Completion}
\label{sec.completion}

Fix $Z,d\in\R_{>0}$ and let $(H_n,\rho_n)_{n\in\N}$ be the directed system
\cref{eq.dyadic_data,eq.dyadic_state_refinement} of isometric embeddings.
Using the increasing-union and completion constructions of
\cref{sec.background_hilbert}, define
\begin{equation}\label{eq.completed_state}
H_\infty(d,Z)
\coloneqq
\bigcup_{n\in\N}H_n,
\qquad
\defineTerm{EZ}(d)\coloneqq\widehat{H_\infty(d,Z)},
\end{equation}
and for duration zero put $\EZ(0)\coloneqq0$.  By
\cref{lem.isometric_colimits}, maps out of $\EZ(d)$ are uniformly
bounded compatible families \cref{eq.completion_universal_property}: it
is the colimit of $(H_n,\rho_n)$ in Hilbert spaces and contractions.

\begin{theorem}[Dyadic completion]\label{thm.dyadic_completion}
Fix $Z,d\in\R_{>0}$.  There is a canonical unitary
\begin{equation}\label{eq.interval_realization}
\iota_Z(d)\colon
\EZ(d)\xrightarrow{\ \cong\ }L^2([0,d),Z\,dx).
\end{equation}
For every $n\in\N$, the image of $H_n$ consists of the functions constant
on each interval
\begin{equation}\label{eq.dyadic_cells}
\left[
\frac{jd}{2^n},
\frac{(j+1)d}{2^n}
\right),
\qquad
0\leq j<2^n.
\end{equation}
\end{theorem}

\begin{proof}
For every $n\in\N$, send
$x=(x_0,\ldots,x_{2^n-1})\in H_n$ to the step function with value $x_j$
on the $j$th interval in \cref{eq.dyadic_cells}.  This map is isometric
because
\[
\sum_{j=0}^{2^n-1}
\int_{jd/2^n}^{(j+1)d/2^n}Zx_j^2\,dx
=
\frac{Zd}{2^n}\sum_{j=0}^{2^n-1}x_j^2
=
\lVert x\rVert_{H_n}^2.
\]
The maps commute with $\rho_n$, so they assemble into an isometry from
$H_\infty(d,Z)$ onto the subspace of dyadic step functions, which is dense
in $L^2([0,d),Z\,dx)$; completion gives \cref{eq.interval_realization}.
\end{proof}

\subsection{Ordered intervals}
\label{sec.intervals}

\begin{definition}[Ordered interval operad]\label{def.interval_operad}
The nonsymmetric colored operad $\defineTerm{Int}$ has colors
$d\in\R_{\geq0}$ and operation sets
\[
\Int(d_1,\ldots,d_k;d)
\coloneqq
\begin{cases}
\{*\},&d=d_1+\cdots+d_k,\\
\varnothing,&d\neq d_1+\cdots+d_k,
\end{cases}
\]
for every integer $k\geq0$ and all
$d,d_1,\ldots,d_k\in\R_{\geq0}$.  Substitution is induced by
associativity of addition.
\end{definition}

\begin{theorem}[The characteristic interval algebra]
\label{thm.characteristic_interval_algebra}
Fix $Z\in\R_{>0}$.  For every integer $k\geq0$ and durations
$d_1,\ldots,d_k\in\R_{\geq0}$, interval concatenation gives a unitary
\[
\Gamma^Z_{d_1,\ldots,d_k}\colon
\bigoplus_{i=1}^k\EZ(d_i)
\xrightarrow{\ \cong\ }
\EZ(d_1+\cdots+d_k).
\]
These maps make
$\EZ\colon\Int\to\Hilb$, $d\mapsto \EZ(d)$,
a strong $\Int$-algebra in $(\Hilb,0,\oplus)$
(\cref{sec.background_categories}).
\end{theorem}

\begin{proof}
Put $s_0\coloneqq0$ and
$s_i\coloneqq d_1+\cdots+d_i$ for $1\leq i\leq k$.  Under
\cref{eq.interval_realization}, the map
$\Gamma^Z_{d_1,\ldots,d_k}$ translates the function on $[0,d_i)$ to
$[s_{i-1},s_i)$.  These intervals are disjoint and cover $[0,s_k)$, so
the map is unitary; for $k=0$ it is the identity of $\EZ(0)=0$.
Translation in stages equals translation all at once,
which is the substitution axiom of $\Int$.
\end{proof}

\begin{definition}[Completed step]\label{def.completed_step}
Fix $Z,d\in\R_{>0}$ and $a\in(0,d]$.  The \emph{completed step} is the
unitary composite
\begin{equation}\label{eq.completed_delay_step}
\step^Z_{d,a}\colon
\EZ(d)\oplus \EZ(a)
\xrightarrow{\ (\Gamma^Z_{a,d-a})^{-1}\oplus\id\ }
\EZ(a)\oplus \EZ(d-a)\oplus \EZ(a)
\xrightarrow{\ \id\oplus\Gamma^Z_{d-a,a}\ }
\EZ(a)\oplus \EZ(d).
\end{equation}
We write $\defineTerm{headf}_{d,a}\colon\EZ(d)\to\EZ(a)$ and
$\defineTerm{nextf}_{d,a}\colon\EZ(d)\oplus\EZ(a)\to\EZ(d)$ for its two
components.
\end{definition}

Under the interval realization \cref{eq.interval_realization}, the
completed step is cut-and-translate: for almost every $x\in[0,d)$,
\begin{equation}\label{eq.cut_translate}
\headf_{d,a}(q)=q|_{[0,a)},
\qquad
\nextf_{d,a}(q,b)(x)
=
\begin{cases}
q(x+a),&0\leq x<d-a,\\
b(x-d+a),&d-a\leq x<d.
\end{cases}
\end{equation}

\begin{lemma}[Block steps are natural]\label{lem.block_naturality}
Fix $Z,d\in\R_{>0}$ and $n\in\N$.  With the data \cref{eq.dyadic_data}
and, for $m\geq n$,
\[
M_{n,m}\coloneqq2^{m-n},
\qquad
B_{n,m}\coloneqq \Qu_{M_{n,m}}(C_m),
\]
the assignments
\[
F_n(m)\coloneqq H_m\oplus B_{n,m},
\qquad
G_n(m)\coloneqq B_{n,m}\oplus H_m,
\]
with transition maps the direct sums of the refinements
$\rho_{K_m,2,\sigma_m}$ and $\rho_{M_{n,m},2,\sigma_m}$
(\cref{thm.queue_subdivision,cor.dyadic_subdivision}), are functors
$(\N_{\geq n},\leq)\to\Hilb$, and the block steps
\[
\defineTerm{blk}_n(m)\coloneqq\step_{K_m}^{[M_{n,m}]}
\colon F_n(m)\longrightarrow G_n(m)
\]
form a natural transformation $\blk_n\colon F_n\Rightarrow G_n$, every
component of which is unitary.
\end{lemma}

\begin{proof}
Each transition map is a direct sum of isometric embeddings
(\cref{cor.dyadic_subdivision}), and a functor from the poset
$(\N_{\geq n},\leq)$ is determined by its consecutive transition maps.
Naturality at $m\leq m+1$ is the queue-subdivision square
\cref{eq.queue_subdivision_square} with $B=C_m$, $C=C_{m+1}$,
$N=K_m$, $L=M_{n,m}$, $r=2$, and $\sigma=\sigma_m$.  Each component is
a composite of queue steps padded by identities, and the queue step is
unitary (\cref{thm.queue_representation}).
\end{proof}

\begin{theorem}[The completed step is the colimit of the block steps]
\label{thm.step_colimit}
Fix $Z,d\in\R_{>0}$ and $n\in\N$, and let
$\jmath_m\colon H_m\to \EZ(d)$ and
$\epsilon_{n,m}\colon B_{n,m}\to \EZ(h_n)$
be the step-function embeddings supplied by
\cref{thm.dyadic_completion}.  In Hilbert spaces and contractions,
\[
\operatorname{colim}F_n=\EZ(d)\oplus \EZ(h_n),
\qquad
\operatorname{colim}G_n=\EZ(h_n)\oplus \EZ(d),
\]
with cocones $\jmath_m\oplus\epsilon_{n,m}$ and
$\epsilon_{n,m}\oplus\jmath_m$, and
\[
\step^Z_{d,h_n}
=
\operatorname{colim}\blk_n:
\]
the completed step is the unique bounded map making the squares
\begin{equation}\label{eq.finite_completed_square}
\begin{CD}
H_m\oplus B_{n,m}
@>{\blk_n(m)}>>
B_{n,m}\oplus H_m\\
@V{\jmath_m\oplus\epsilon_{n,m}}VV
@VV{\epsilon_{n,m}\oplus\jmath_m}V\\
\EZ(d)\oplus \EZ(h_n)
@>>{\step^Z_{d,h_n}}>
\EZ(h_n)\oplus \EZ(d)
\end{CD}
\end{equation}
commute for every $m\geq n$.
\end{theorem}

\begin{proof}
Direct sums commute with increasing unions and with Hilbert
completion, so the completed union of $F_n$ is
$\EZ(d)\oplus\EZ(h_n)$: apply \cref{thm.dyadic_completion} to duration
$d$ for the first summand and to duration $h_n$ for the second, whose
dyadic tower at resolutions $m\geq n$ is exactly $(B_{n,m})$, since
$C_m=\CZ(h_n/2^{m-n})$.  Likewise for $G_n$.  By
\cref{lem.isometric_colimits}, these completed unions are the colimits
in Hilbert spaces and contractions, with the stated cocones.

The composites
$(\epsilon_{n,m}\oplus\jmath_m)\circ\blk_n(m)$ are isometries, hence
uniformly bounded, and they are a compatible family: naturality
(\cref{lem.block_naturality}) exchanges $\blk_n$ with the transition
maps, and the embeddings commute with the transition maps
(\cref{thm.dyadic_completion}).  By \cref{lem.isometric_colimits}
exactly one bounded map out of $\operatorname{colim}F_n$ restricts to
this family; that map is $\operatorname{colim}\blk_n$, and
\cref{eq.finite_completed_square} is its defining square.  It remains
to identify it with $\step^Z_{d,h_n}$ \cref{eq.completed_delay_step}:
on the dyadic step functions both maps cut off the first interval of
duration $h_n$, record it as output, and append the input
interval---the block step packet by packet, the completed step by
\cref{eq.cut_translate}.
\end{proof}

\section{The divisible string}
\label{sec.divisible}

\begin{definition}[Divisible string]\label{def.divisible_string}
Fix $Z\in\R_{>0}$, a tick duration $a\in\R_{>0}$, and $d\in\R_{\geq a}$.
The \emph{divisible string} is the ported machine
\[
\Sm_Z(d;a)\colon\bigl(\EZ(a),0\bigr)\longrightarrow\bigl(\EZ(a),0\bigr)
\quad\text{in }\Lin
\]
with state space $\EZ(d)\oplus\EZ(d)$ and
\[
\begin{gathered}
\Sm_Z(d;a)^+(x,y)
\coloneqq
\Bigl(\headf_{d,a}(x),\
\riesz_{\EZ(a)}^{-1}\headf_{d,a}(y)\Bigr),\\
\Sm_Z(d;a)^-\bigl((x,y),\alpha,\theta\bigr)
\coloneqq
\Bigl(\nextf_{d,a}(x,\alpha),\
\nextf_{d,a}\bigl(y,\riesz_{\EZ(a)}(\theta)\bigr)\Bigr).
\end{gathered}
\]
\end{definition}

The rightward channel $x$ carries value packets toward the output
port; the leftward channel $y$ carries covector packets, transported
through the Riesz unitaries of $\EZ(a)$, toward the input port.  The
finite strings of \cref{def.string} are the dyadic points of this
definition: \cref{thm.string_colimit} exhibits $\Sm_Z(d;h_n)$ as
their colimit.

\begin{theorem}[Divisibility]\label{thm.divisibility}
Fix $Z\in\R_{>0}$ and a tick duration $a\in\R_{>0}$.  For
$d_1,d_2\in\R_{\geq a}$ and the interval concatenations $\Gamma^Z$ of
\cref{thm.characteristic_interval_algebra}, the map
\[
\chi^Z_{d_1,d_2}\bigl((x,y),(x',y')\bigr)
\coloneqq
\Bigl(
\bigl(\Gamma^Z_{d_2,d_1}(x',x),\
\Gamma^Z_{d_1,d_2}(y,y')\bigr)
\Bigr)
\]
is a unitary $2$-cell
\[
\chi^Z_{d_1,d_2}\colon
\Sm_Z(d_2;a)\circ\Sm_Z(d_1;a)
\xrightarrow{\ \cong\ }
\Sm_Z(d_1+d_2;a),
\]
and the family $(\chi^Z_{d_1,d_2})$ satisfies the associativity
coherence: the assignment $d\mapsto\Sm_Z(d;a)$ is a strong semigroupal
representation of $(\R_{\geq a},+)$ in
$\bigl(\Lin((\EZ(a),0),(\EZ(a),0)),\circ\bigr)$.
\end{theorem}

\begin{proof}
Each concatenation $\Gamma^Z$ is unitary
(\cref{thm.characteristic_interval_algebra}).  We verify the machine
commutations for the rightward channel in the function picture
\cref{eq.cut_translate}; the leftward channel is the mirror
computation, conjugated by the Riesz unitaries.  Write
$x''\coloneqq\Gamma^Z_{d_2,d_1}(x',x)$, the function equal to $x'$ on
$[0,d_2)$ and to $x(\blank-d_2)$ on $[d_2,d_1{+}d_2)$; the exit port
of the composite is that of the second string, so the near-exit
content is $x'$.

Readout: the composite's value \cref{eq.machine_composition} is the
second string's, $\headf_{d_2,a}(x')=x'|_{[0,a)}=x''|_{[0,a)}
=\headf_{d_1+d_2,a}(x'')$.

Update: the middle impedance form is zero, so the first string
receives the input $\alpha$ and the second receives the first's value
$\headf_{d_1,a}(x)$.  The updated concatenation is the function
\[
s\longmapsto
\begin{cases}
x'(s+a),&s<d_2-a,\\
x(s+a-d_2),&d_2-a\leq s<d_2,\\
x(s+a-d_2),&d_2\leq s<d_1+d_2-a,\\
\alpha(s+a-d_1-d_2),&d_1+d_2-a\leq s,
\end{cases}
\]
where the first two cases are $\nextf_{d_2,a}(x',
\headf_{d_1,a}x)$ and the last two are
$\nextf_{d_1,a}(x,\alpha)$ placed on $[d_2,d_1{+}d_2)$; all four
agree with $\nextf_{d_1+d_2,a}(x'',\alpha)(s)$ by
\cref{eq.cut_translate}, since $x''(s+a)=x'(s+a)$ for $s+a<d_2$ and
$x''(s+a)=x(s+a-d_2)$ for $s+a\geq d_2$.

Coherence: on a triple composite, both composites of concatenations
are the single concatenation reordered by the substitution axiom of
$\Int$ (\cref{thm.characteristic_interval_algebra}), i.e.\ the same
unitary.
\end{proof}

\begin{theorem}[The divisible string is the colimit of the finite strings]
\label{thm.string_colimit}
Fix $Z,d\in\R_{>0}$ and $n\in\N$, with the data \cref{eq.dyadic_data}
and $B_{n,m}$ as in \cref{lem.block_naturality}.

\begin{enumerate}
\item\label{item.sc_colim} The packet refinements
$\rho_{M_{n,m},2,\sigma_m}$ make $m\mapsto\Sm_{B_{n,m}}(K_n)$ a tower
of ported machines, and
\[
\operatorname{colim}_m\Sm_{B_{n,m}}(K_n)
\cong
\Sm_{\EZ(h_n)}(K_n).
\]

\item\label{item.sc_ident} For every $a\in\R_{>0}$ and
$K\in\Z_{\geq1}$, writing
$\Gamma^{(K)}_a\coloneqq\Gamma^Z_{a,\ldots,a}$ for the $K$-fold
interval concatenation
(\cref{thm.characteristic_interval_algebra}), the map
\[
\iota_{K,a}
\coloneqq
\Gamma^{(K)}_a
\oplus
\Bigl(\Gamma^{(K)}_a\circ \Qu_K(\riesz_{\EZ(a)})\Bigr)
\colon
\Sm_{\EZ(a)}(K)\To{\cong}\Sm_Z(Ka;a)
\]
is a unitary $2$-cell isomorphism in
$\Lin\bigl((\EZ(a),0),(\EZ(a),0)\bigr)$.  In particular
$\Sm_Z(d;h_n)\cong\operatorname{colim}_m\Sm_{B_{n,m}}(K_n)$.

\item\label{item.sc_chi} The isomorphisms of \cref{item.sc_ident}
carry the concatenation $2$-cells of \cref{prop.string_composition}
to those of \cref{thm.divisibility}: for all $K_1,K_2\in\Z_{\geq1}$,
\[
\chi^Z_{K_1a,K_2a}
\circ\bigl(\iota_{K_1,a}\oplus\iota_{K_2,a}\bigr)
=
\iota_{K_1+K_2,a}\circ\chi_{K_1,K_2}.
\]
\end{enumerate}
\end{theorem}

\begin{proof}
\Cref{item.sc_colim}: the packet refinements are isometries
(\cref{cor.dyadic_subdivision}) between zero-impedance ports, so
\cref{prop.string_functorial} produces the tower, and its state maps
are isometries with contractive structure maps, so
\cref{prop.machine_colimits} applies.  The completed union of the
alphabet tower $(B_{n,m})_{m\geq n}$ is $\EZ(h_n)$ by
\cref{thm.dyadic_completion}, since $C_m=\CZ(h_n/2^{m-n})$, and the
string functor preserves the colimit
(\cref{prop.string_functorial}).

\Cref{item.sc_ident}: $\iota_{K,a}$ is unitary, being assembled from
the concatenations $\Gamma^{(K)}_a$
(\cref{thm.characteristic_interval_algebra}) and the Riesz unitaries.
Under the interval realization, $\Gamma^{(K)}_a$ places slot $j$ on
$[ja,(j+1)a)$, so by \cref{eq.cut_translate}
\[
\headf_{Ka,a}\circ\Gamma^{(K)}_a=P_K
\qquad\text{and}\qquad
\nextf_{Ka,a}\circ\bigl(\Gamma^{(K)}_a\oplus\id\bigr)
=\Gamma^{(K)}_a\circ U_K
\]
for the queue maps of \cref{eq.queue_structure}: cutting off the
first interval reads the head slot, and translating with an appended
input shifts with an inserted input.  These are the value-channel
commutations of a $2$-cell (\cref{def.string,def.divisible_string}),
and the covector-channel commutations are their conjugates by the
slotwise Riesz unitaries.  The last claim combines
\cref{item.sc_colim,item.sc_ident} at $a=h_n$ and $K=K_n$, where
$K_nh_n=d$.

\Cref{item.sc_chi}: on the rightward channel, the two sides applied
to $\bigl((x,y),(x',y')\bigr)$ are
$\Gamma^Z_{K_2a,K_1a}\circ
\bigl(\Gamma^{(K_2)}_a\oplus\Gamma^{(K_1)}_a\bigr)$ and
$\Gamma^{(K_1+K_2)}_a$, both evaluated on the $K_1{+}K_2$ slots of
$(x',x)$: concatenation in stages equals concatenation all at once,
the substitution axiom of $\Int$
(\cref{thm.characteristic_interval_algebra}).  The leftward channel
is the same identity applied to $(y,y')$, conjugated by the slotwise
Riesz unitaries.
\end{proof}

At commensurable splits, divisibility is therefore the finite
concatenation isomorphism transported along \cref{item.sc_ident}; the
splits of \cref{thm.divisibility} at arbitrary real durations are
supplied by the interval algebra.

\subsection{The main theorem}
\label{sec.main}

\begin{theorem}[The divisible string is the infinite-resolution chain]
\label{thm.main}
Fix $Z,d\in\R_{>0}$, and for $n\in\N$ put $K_n\coloneqq2^n$,
$h_n\coloneqq d/2^n$, $C_n\coloneqq\CZ(h_n)$, and
$H_n\coloneqq\Qu_{K_n}(C_n)$ with queue shift $S_{K_n}$
\cref{eq.queue_structure}.

\begin{enumerate}
\item\label{item.main_syntax} \textup{(Chains compose.)}  The
assignment $K\mapsto\Chain_K$ is a homomorphism of semigroups
$(\Z_{\geq1},+)\to\Open(\Zport,\Zport)$, and the open linear semantics
carries it to the chain machines:
$\Tm_Z(K{+}K')=\Tm_Z(K')\circ\Tm_Z(K)$ in $\Lin$.

\item\label{item.main_flow} \textup{(At every resolution, the chain
is a string in wave variables.)}  For every driven trajectory of
$\Tm_Z(K_n)$ and every $t\geq K_n$, the wave chart
\cref{def.packets} equals the state of the string $\Sm_{C_n}(K_n)$
driven by the in-streams \cref{eq.in_streams}, and the out-streams of
the two machines coincide: they are the in-streams delayed by exactly
$K_n$ ticks.

\item\label{item.main_completion} \textup{(The strings complete to
cut-and-translate.)}  Packet doubling and time change along
$2\cdot\colon\N\to\N$ embed each resolution into the next,
$\rho_n\colon(H_n,S_{K_n})\to(2\cdot)^*(H_{n+1},S_{K_{n+1}})$ in
$\Dyn$; the completed union of either travel direction of the tower is
$\EZ(d)\cong L^2\bigl([0,d),Z\,dx\bigr)$; the completed steps are
the colimits of the natural transformations of block steps,
$\step^Z_{d,h_n}=\operatorname{colim}\blk_n$, which are
cut-and-translate; and at the level of machines,
$\Sm_Z(d;h_n)\cong\operatorname{colim}_m\Sm_{B_{n,m}}(K_n)$, the
colimit of the finite strings over the block alphabets.

\item\label{item.main_divisibility} \textup{(Divisibility.)}  For
every tick duration $a\in\R_{>0}$, the assignment
$d'\mapsto\Sm_Z(d';a)$ is a strong semigroupal representation of
$(\R_{\geq a},+)$ in
$\bigl(\Lin((\EZ(a),0),(\EZ(a),0)),\circ\bigr)$: concatenation of
intervals is serial composition of strings.
\end{enumerate}
\end{theorem}

\begin{proof}
\Cref{item.main_syntax} is
\cref{cor.chain_semigroup,cor.chain_machines_compose}.
\Cref{item.main_flow} is
\cref{thm.characteristic_flow,cor.chain_string}: the structure maps
\cref{eq.queue_structure} of $\Sm_B(K)$ depend only on the underlying
vector space of $B$, so the comparison over $\R$ is the comparison
over the weighted packet space $C_n$.
\Cref{item.main_completion} is
\cref{cor.dyadic_subdivision,thm.dyadic_completion,lem.block_naturality,thm.step_colimit,thm.string_colimit},
with the leftward tower conjugated by the slotwise Riesz unitaries of
the $C_n$.
\Cref{item.main_divisibility} is \cref{thm.divisibility}.
\end{proof}

The infinite-resolution chain is thus reached from the cell in four
typed steps: powers of the cell in $\Open$
\cref{eq.open_composition}, the semantics $\Phiopen$
(\cref{thm.open_semantics}), the characteristic flow
\cref{eq.transport}, and the completion
\cref{eq.completion_universal_property}.  Every step preserves
composition: gluing of cells, serial composition of machines,
concatenation of queues, concatenation of intervals.

\section{Continuous waves}
\label{sec.waves}

This section uses local $L^2$ analysis that the rest of the paper does
not, so we record that background here rather than in
\cref{sec.background_hilbert}.

\paragraph{Local spaces and distributions.}
Let $\Omega\subseteq\R^2$ be open.  We write $L^2_{\mathrm{loc}}(\Omega)$
for the space of measurable functions square-integrable on every compact
subset.  Every $u\in L^2_{\mathrm{loc}}(\Omega)$ is a distribution on
$\Omega$ and has distributional partial derivatives; we write
$H^1_{\mathrm{loc}}(\Omega)$ for the space of
$u\in L^2_{\mathrm{loc}}(\Omega)$ whose first distributional derivatives
again lie in $L^2_{\mathrm{loc}}(\Omega)$
\cite{brezis2011functional}.  For a bounded interval $J\subseteq\R$, the
Sobolev space $H^1(J)$ consists of the $u\in L^2(J)$ with distributional
derivative $\partial_xu\in L^2(J)$; we write $H^1(J)/\R$ for its quotient
by the constant functions.  Translation acts strongly continuously on
$L^2_{\mathrm{loc}}$: for $u\in L^2_{\mathrm{loc}}$, the translates of $u$
converge to $u$ in $L^2$ on compacta as the displacement tends to $0$.  We
use the distributional Poincar\'e lemma on rectangles: if
$v,e\in L^2_{\mathrm{loc}}(\Omega)$ with $\Omega$ an open rectangle satisfy
$\partial_\tau e=\partial_xv$ distributionally, then the closed
distributional one-form
$v\,d\tau+e\,dx$ has a potential $u\in H^1_{\mathrm{loc}}(\Omega)$---that
is, $\partial_\tau u=v$ and $\partial_xu=e$---unique
up to an additive constant \cite{schwartz1966distributions}.

\begin{theorem}[Characteristic-wave correspondence]
\label{thm.characteristic_wave}
Let $\Omega\subseteq\R^2$ be an open rectangle with coordinates
$(\tau,x)$, and let
$w^{\leftarrow},w^{\rightarrow}\in L^2_{\mathrm{loc}}(\Omega)$.
The following are equivalent.

\begin{enumerate}
\item For every $a\in\R_{>0}$ for which the displayed arguments belong to
$\Omega$, the local translation laws
\begin{equation}\label{eq.local_translation_laws}
w^{\leftarrow}(\tau+a,x)
=
w^{\leftarrow}(\tau,x+a),
\qquad
w^{\rightarrow}(\tau+a,x)
=
w^{\rightarrow}(\tau,x-a)
\end{equation}
hold for almost every $(\tau,x)\in\Omega$.

\item The distributional transport equations
\begin{equation}\label{eq.transport_equations}
(\partial_\tau-\partial_x)w^{\leftarrow}=0,
\qquad
(\partial_\tau+\partial_x)w^{\rightarrow}=0
\end{equation}
hold on $\Omega$.

\item There is a function $u\in H^1_{\mathrm{loc}}(\Omega)$, unique up to
an additive constant, satisfying
\begin{equation}\label{eq.characteristic_reconstruction}
\partial_\tau u
=
w^{\leftarrow}+w^{\rightarrow},
\qquad
\partial_xu
=
w^{\leftarrow}-w^{\rightarrow},
\end{equation}
and the weak wave equation
\begin{equation}\label{eq.wave_equation}
\partial_\tau^2u=\partial_x^2u.
\end{equation}
\end{enumerate}

Conversely, for every
$u\in H^1_{\mathrm{loc}}(\Omega)$ satisfying
\cref{eq.wave_equation}, the formulas
\begin{equation}\label{eq.wave_characteristics}
w^{\leftarrow}
\coloneqq
\frac{\partial_\tau u+\partial_xu}{2},
\qquad
w^{\rightarrow}
\coloneqq
\frac{\partial_\tau u-\partial_xu}{2}
\end{equation}
define the unique characteristic fields satisfying
\cref{eq.characteristic_reconstruction}.
\end{theorem}

\begin{proof}
Strong continuity of translation on $L^2_{\mathrm{loc}}(\Omega)$
(recorded at the start of this section) identifies
the difference quotients of \cref{eq.local_translation_laws} with the two
distributional derivatives in \cref{eq.transport_equations}.  Conversely,
the coordinates $(\tau,x+\tau)$ for $w^{\leftarrow}$ and
$(\tau,x-\tau)$ for $w^{\rightarrow}$ show that
\cref{eq.transport_equations} makes each field constant along its
characteristics, which gives \cref{eq.local_translation_laws}.

Put
$v\coloneqq w^{\leftarrow}+w^{\rightarrow}$ and
$e\coloneqq w^{\leftarrow}-w^{\rightarrow}$.  The transport equations give
\begin{equation}\label{eq.closed_energy_form}
\partial_\tau e=\partial_xv.
\end{equation}
Hence the distributional one-form $v\,d\tau+e\,dx$ is closed.  Since
$\Omega$ is an open rectangle, the distributional Poincar\'e lemma
(recorded at the start of this section) gives
$u\in H^1_{\mathrm{loc}}(\Omega)$ satisfying
\cref{eq.characteristic_reconstruction}, uniquely up to a constant.
Equation \cref{eq.closed_energy_form} then gives
\cref{eq.wave_equation}.  The converse follows by substituting
\cref{eq.wave_characteristics} into \cref{eq.transport_equations}.
\end{proof}

\begin{corollary}[Energy coordinates]\label{cor.energy_coordinates}
Fix $Z,d\in\R_{>0}$ and let $J\subseteq\R$ be an interval of duration $d$.
Define
\[
\mathcal E_Z(J)
\coloneqq
\bigl(H^1(J)/\R\bigr)\oplus L^2(J)
\]
with squared norm
\[
\lVert([\varphi],v)\rVert_{\mathcal E_Z(J)}^2
\coloneqq
\frac Z2\int_J
\bigl(|\partial_x\varphi|^2+|v|^2\bigr)\,dx.
\]
After translating $J$ to $[0,d)$, the map
\[
\kappa_{Z,J}\colon
\mathcal E_Z(J)\xrightarrow{\ \cong\ }\EZ(d)\oplus \EZ(d),
\qquad
\kappa_{Z,J}([\varphi],v)
\coloneqq
\left(
\frac{v-\partial_x\varphi}{2},
\frac{v+\partial_x\varphi}{2}
\right)
\]
is unitary.
\end{corollary}

\begin{proof}
For every $([\varphi],v)\in\mathcal E_Z(J)$,
\[
\left|\frac{v-\partial_x\varphi}{2}\right|^2
+
\left|\frac{v+\partial_x\varphi}{2}\right|^2
=
\frac12\bigl(|v|^2+|\partial_x\varphi|^2\bigr),
\]
so $\kappa_{Z,J}$ is isometric.  Given
$f,g\in L^2(J,Z\,dx)$, put $v=f+g$ and $e=g-f$.  There is a unique
$[\varphi]\in H^1(J)/\R$ with $\partial_x\varphi=e$, proving surjectivity.
\end{proof}

A trajectory of the divisible string realizes the equivalence of
\cref{thm.characteristic_wave} concretely.  Transporting the two
channels of $\Sm_Z(d;a)$ along their characteristics---the rightward
channel by decreasing time-to-exit, the leftward channel oppositely,
with the in-streams filling in behind them---defines a pair of fields
on the strip $\R_{>0}\times(0,d)$ satisfying the translation laws
\cref{eq.local_translation_laws} by construction, whose time slices
are the states of the trajectory.  By
\cref{thm.characteristic_wave}, these fields are the characteristic
pair \cref{eq.wave_characteristics} of a weak solution of the wave
equation, unique up to an additive constant, and by
\cref{cor.energy_coordinates} each state of the string is a point of
that solution's energy space $\mathcal E_Z\bigl([0,d)\bigr)$.  The
displacement that the wave adds to the two travelling fields is
exactly the constant mode that the quotient $H^1(J)/\R$ removes.

\begin{thebibliography}{9}

\bibitem{cdlm}
D.~I. Spivak.
\newblock \emph{Compositional Dynamics in Learning and Mechanics}.
\newblock Draft, 2026.  Cited by the numbering of the July 2026 draft.

\bibitem{niu2025polynomial}
N.~Niu and D.~I. Spivak.
\newblock \emph{Polynomial Functors: A Mathematical Theory of Interaction}.
\newblock London Mathematical Society Lecture Note Series.  Cambridge
University Press, 2025.

\bibitem{conway1990functional}
J.~B. Conway.
\newblock \emph{A Course in Functional Analysis}.
\newblock Second edition, Graduate Texts in Mathematics 96.  Springer, 1990.

\bibitem{heunen2019categories}
C.~Heunen and J.~Vicary.
\newblock \emph{Categories for Quantum Theory: An Introduction}.
\newblock Oxford University Press, 2019.

\bibitem{cartan1971differential}
H.~Cartan.
\newblock \emph{Differential Calculus}.
\newblock Hermann, 1971.

\bibitem{brezis2011functional}
H.~Brezis.
\newblock \emph{Functional Analysis, Sobolev Spaces and Partial Differential
Equations}.
\newblock Springer, 2011.

\bibitem{schwartz1966distributions}
L.~Schwartz.
\newblock \emph{Th\'eorie des distributions}.
\newblock Hermann, 1966.

\bibitem{leinster2004higher}
T.~Leinster.
\newblock \emph{Higher Operads, Higher Categories}.
\newblock London Mathematical Society Lecture Note Series 298.  Cambridge
University Press, 2004.

\end{thebibliography}

\end{document}
